% !TeX root = ncku_thesis.tex
% this file is encoded in utf-8
% v2.02 (Sep. 12, 2012)

\documentclass[12pt, a4paper]{ncku_report}

% To have correct unicode mapping for special LaTeX characters (italic, ligature, ...) to unicode 
% in the generated PDF file, so that we can do copy-and-paste.
% This only works for pdflatex work flow; for latex + dvipdfmx work flow, this does not do anything
\usepackage{cmap}
%\usepackage[T1]{fontenc}


% 除非校方修改了論文格式 (margins, header, footer, 浮水印, 中文數字之章別)
% 或者需要增加所用的 LaTeX 套件，
% 或者要改預設中文字型、編碼
% 否則毋須修改本檔內容
% 論文撰寫，請修改以 my_  開頭檔名的各檔案
\usepackage{CJKutf8}  %%% ZZZ %%% macro for Chinese/Japanese/Korean processing
\usepackage{CJKnumb} %%% ZZZ %%% for Chinese numbering capability
\usepackage{CJKspace} %%% ZZZ %%% suppress the spaces due to hard-coded CRs in the text, but not the spaces between CJK chars and alphabets
%\usepackage[nospace]{cite}  % for smart citation
\usepackage{geometry}  % for easy margin settings
%
% margins setting
\geometry{verbose,a4paper,tmargin=3.5cm,bmargin=2cm,lmargin=2cm,rmargin=2cm}
%

% use colored text, background, foreground
% http://en.wikibooks.org/wiki/LaTeX/Colors
\usepackage[usenames,dvipsnames,svgnames,table]{xcolor}

% 插圖套件 graphicx
% 使用者工作流程是用 pdftex 還是 latex + dvipdfmx?
% 視情況而有不同的參數
% 這裡作自動判斷
% 參考自
% http://www.tex.ac.uk/cgi-bin/texfaq2html?label=ifpdf
\newcommand\mydvipdfmxflow{dvipdfmx}
\newcommand\mypdftexflow{pdftex}
\ifx\pdfoutput\undefined
  % not running pdftex
  \usepackage[dvipdfmx]{graphicx}
  \newcommand\myworkflow{dvipdfmx}  % set the flag for hyperref
\else
  \ifx\pdfoutput\relax
    % not running pdftex
    \usepackage[dvipdfmx]{graphicx}
    \newcommand\myworkflow{dvipdfmx}  % set the flag
  \else
    % running pdftex, with...
    \ifnum\pdfoutput>0
      % ... PDF output
      \usepackage[pdftex]{graphicx}
      \newcommand\myworkflow{pdftex}  % set the flag
    \else
      %...DVI output
      \usepackage[dvipdfmx]{graphicx}
      \newcommand\myworkflow{dvipdfmx}  % set the flag
    \fi
  \fi
\fi

% 增強功能型頁楣 / 頁腳套件
\usepackage{fancyhdr}  % 借用此套件來擺放浮水印 
% (佔用了 central header)
% 不需要浮水印的使用者仍可利用此套件，產生所需的 header, footer
%
% 啟動 fancy header/footer 套件
\pagestyle{fancy}
\fancyhead{}  % reset left, central, right header to empty
\fancyfoot[C]{\thepage} %中間 footer 擺放頁碼
\renewcommand{\headrulewidth}{0pt} % header 的直線; 0pt 則無線

% 如果不需要任何浮水印，則請把下列介於 >>> 與 <<< 之間
% 的文字行關掉 (行首加上百分號)
%% 浮水印 >>> 
%
% this file is encoded in utf-8
% v2.02 (Sep. 12, 2012)
% 如果浮水印不是全篇需要，請把下列介於 >>> 與 <<<
% 的「全篇浮水印專用碼」關掉 (行首加百分號)
% 參考自 Keith Reckdahl 寫的 "Using Imported Graphics in LATEX2e" (epslatex.pdf) p.39
% 如果只有特定頁需要浮水印
% 則依該頁屬性使用下列之一的命令 
% 普通頁命令 \thispagestyle{WaterMarkPage}
% plain 頁命令 \thispagestyle{PlainWaterMarkPage}
% empty 頁命令 \thispagestyle{EmptyWaterMarkPage}


% 將重複使用的浮水印章
% 圖檔是 my_watermark.xxx
% 副檔名可以不加，可以是 latex 系統能處裡的任何格式：pdf, gif, png, jpg, eps, ...
% 某些圖檔格式在某些工作流程可能需要作前置處裡。
% 例如，pdflatex 無法直接處理 eps 檔
%  latex + dvipdfmx 無法直接處理 pdf, gif, png, jpg, 需要用 ebb 小工具程式
%  對圖檔產生 .bb 對應檔。
%
% 寬為 5.1 cm
\newsavebox{\mywatermark}
%\sbox{\mywatermark}{\includegraphics[keepaspectratio,width=5.1cm]{image/ncku_watermark.jpg}}
\sbox{\mywatermark}{\includegraphics[keepaspectratio]{image/ncku_watermark.jpg}}

% 將 central header 擺放浮水印的巨集指令
\newcommand{\PlaceWaterMark}{\fancyhead[C]{\setlength{\unitlength}{1in}%
\begin{picture}(0,0)%
\put(-1.5,-6.0){\usebox{\mywatermark}}% 圖檔擺放的位置座標
\end{picture}}%
}

\fancyhead{}  % reset left, central, right header to empty
%% 如果不需整篇論文都要浮水印
%% 則下面  >>> 與 <<< 之間的程式碼請關閉
%% >>> 全篇浮水印
\PlaceWaterMark  % 每一頁都有浮水印 (除了 plain、empty 頁以外)

% 重新定義 plain 頁面
% 每張 plain 頁面 (每一章的第一頁) 也加浮水印

\fancypagestyle{plain}{%
\fancyhead{}%
\PlaceWaterMark%
\fancyfoot{}%
\fancyfoot[C]{\thepage}
\renewcommand{\headrulewidth}{0pt}%
\renewcommand{\footrulewidth}{0pt}%
}
%% <<< 全篇浮水印

%% 如果只有一、兩頁需要有浮水印
%% 可以在該頁 (有頁碼) 使用 \thispagestyle{WaterMarkPage}
%% 此命令不影響原有的 header、footer
\fancypagestyle{WaterMarkPage}{%
\PlaceWaterMark%
}

%% 如果只有一、兩頁 plain 頁需要有浮水印 (如 摘要、自傳等)
%% 可以在該頁 (有頁碼) 使用 \thispagestyle{PlainWaterMarkPage}
%% 只有頁碼與浮水印，沒有其他的 header、footer
%% 等同於 plain page style + water mark
\fancypagestyle{PlainWaterMarkPage}{%
\fancyhead{}%
\PlaceWaterMark%
\fancyfoot{}%
\fancyfoot[C]{\thepage}
\renewcommand{\headrulewidth}{0pt}%
\renewcommand{\footrulewidth}{0pt}%
}

%% 如果只有一、兩頁 empty 頁需要有浮水印 (如封面、書名頁)
%% 可以在該頁 (無頁碼) 使用 \thispagestyle{EmptyWaterMarkPage}
%% 等同於 empty page style + water mark
\fancypagestyle{EmptyWaterMarkPage}{%
\fancyhead{}%
\PlaceWaterMark%
\fancyfoot{}%
\renewcommand{\headrulewidth}{0pt}%
\renewcommand{\footrulewidth}{0pt}%
}

%% <<< 浮水印

% 如需額外的頁楣 (header) 或 footer，請在 my_headerfooter.tex 裡依例修改
% 它的預設內容是都關掉，可依需要打開
%
% this file is encoded in utf-8
% v2.02 (Sep. 12, 2012)

%%%%%%% 其他的 header (left, right) 定義
% 底下定義了一些常見的 header 型式
% 預設情況是關掉的
% 使用者可以視需要將之打開
% 也就是把下列介於 >>> 與 <<< 之間
% 的文字行打開 (行首去掉百分號)

%% header >>>
%\renewcommand{\chaptermark}[1]{%
%\markboth{\prechaptername\ \thechapter\ \postchaptername%
%\ #1}{}%
%}  %定義 header 使用的「章」層級的戳記
%\fancyhead[L]{} % 左 header 為空
%\fancyhead[R]{\leftmark}  % 右 header 擺放「章」層級的戳記 (以 \leftmark 叫出)
%\renewcommand{\headrulewidth}{0.4pt}  % header 的直線 0.4pt; 0pt 則無線
%% <<< header

%%%%%%% 其他的 footer (left, right) 定義
% 底下定義了一些常見的 footer 型式
% 預設情況是關掉的
% 使用者可以視需要將之打開
% 也就是把下列介於 >>> 與 <<< 之間
% 的文字行打開 (行首去掉百分號)

%% footer >>>
%\fancyfoot[L]{} % 左 footer 為空
%\fancyfoot[R]{\small{YZU \LaTeX\ v2.0}} % 右 footer 擺放論文格式版本
%\renewcommand{\footrulewidth}{0.4 pt} % footer 的直線 0.4pt; 0pt 則無線
%% <<< footer




%%%%%%%%%%%%%%%%%%%%%%%%%%%%%%
%%%% 非必要的套件，但很實用
\usepackage{amsmath} % 各式 AMS 數學功能
\usepackage{amssymb} % 各式 AMS 數學符號
\usepackage{mathrsfs} %草寫體數學符號，在數學模式裡用 \mathscr{E} 得草寫 E
\usepackage{listings} % 程式列表套件
\usepackage{booktabs} % For better table formatting
\usepackage{graphicx} % Required for including images
\usepackage{array, float}    % For better columns
\usepackage{multirow} % Add this to use the multirow functionality
\usepackage{tikz}
\usepackage{caption}
\usepackage{subcaption}

%
% listing setting
\lstset{breaklines=true,% 過長的程式行可斷行
extendedchars=false,% 中文處理不需要 extendedchars
texcl=true,% 中文註解需要有 TeX 處理過的 comment line, 所以設成 true
comment=[l]\%\%,% 以雙「百分號」做為程式中文註解的起頭標記，配合 MATLAB
basicstyle=\small,% 小號字體, 約 10 pt 大小
commentstyle=\upshape,% 預設是斜體字，會影響註解裏的英文，改用正體
escapeinside={<>},%要在列表裡顯示中文字串，要把該中文字串用 <> 包夾住
%language=Octave % 會將一些 octave 指令以粗體顯示
}

\usepackage{url} % 在文稿中引用網址，可以用 \url{http://www.ncku.edu.tw} 方式

\usepackage{CJKfntef}  %%% ZZZ %%% CJK underline: \CJKunderline{}
\normalem %%% ZZZ %%% use \normalem to avoid \emph becomes underline, due to CJKfntef

% the following packages are needed for table code generated by excel2latex macro in excel
\usepackage{multirow} % for sophisticated table with columns spanning multiple rows; \multirow{num_rows}{width}{contents}
\usepackage{bigstrut} % to have larger cell area
\usepackage{booktabs} % to have a more professional look of tables

% use colored text, background, foreground
% http://en.wikibooks.org/wiki/LaTeX/Colors
%\usepackage[usenames,dvipsnames,svgnames,table]{xcolor}

%didi add
\usepackage[vlined, ruled, linesnumbered]{algorithm2e}
\usepackage[noend]{algpseudocode}
\usepackage{amsmath}
\let\proof\rela
\let\endproof\relax
\usepackage{amsthm}
\newtheorem{theorem}{Theorem}
\newtheorem{lemma}{Lemma}
\newtheorem{definition}{Definition}
\renewcommand{\algorithmicrequire}{\textbf{Input:}}  % Use Input in the format of Algorithm
\renewcommand{\algorithmicensure}{\textbf{Output:}} % Use Output in the format of Algorithm
\usepackage{diagbox}
\usepackage{amssymb}
\usepackage{multirow}
\usepackage{amsfonts}
\usepackage{thmtools}
\newcommand{\tabincell}[2]{\begin{tabular}{@{}#1@{}}#2\end{tabular}}  
\usepackage{graphicx}
\usepackage{url}
%didi add

% Add biblatex for bibliography management
\usepackage[
  backend=bibtex,
  style=numeric,
  sorting=none
]{biblatex}
\addbibresource{texFolder/reference.bib}

%%%% 以上為非必要套件
%%%%%%%%%%%%%%%%%%%%%%%%%%%%%%

%%% 以下是 hyperref 套件
%%%%%%%%%%%%%%%%%%%%%%%%%%%%%%
% hyperref 會擾亂 cite.sty 對文獻號碼縮編的排版，所以依據
% http://www.ctan.org/tex-archive/macros/latex/contrib/hyperref/
% 作如下的更動，使得 hyperref 不做文獻號碼的超連結。
% \makeatletter
% \def\NAT@parse{\typeout{This is a fake Natbib command to fool Hyperref.}}
% \makeatother

% hyperlinkable table of contents
% 章節目錄、圖表超連結
\ifx\myworkflow\mydvipdfmxflow
        \usepackage[dvipdfmx, debug, colorlinks, linkcolor=black, citecolor=black, urlcolor=black, unicode]{hyperref}
\else
        \usepackage[pdftex, debug, colorlinks, linkcolor=black, citecolor=black, urlcolor=black, unicode]{hyperref}      
\fi

% if hyperref is not used (e.g., in LyX application)
% define dummy \phantomsection for those occurences
%   in ncku_frontpages.tex, ncku_backpages.tex, my_appendix.tex
\ifx\hypersetup\undefined
        \newcommand\phantomsection{}
\fi
%%%% 以上為所有套件
%%%% 
%%%% 


% global page layout
\newcommand{\mybaselinestretch}{1.5}  %行距 1.5 倍 + 20%, (約為 double space)
\renewcommand{\baselinestretch}{\mybaselinestretch}  % 論文行距預設值
\parskip=2ex  % 段落之間的間隔為兩個 x 的高度
\parindent = 0Pt  % 段首內縮由 CJK 控制，所以這裡就設成不內縮

%%%%%%%%%%%%%%%%%%%%%%%%%%%%%
%  end of preamble
%%%%%%%%%%%%%%%%%%%%%%%%%%%%%
%
\usepackage{siunitx}
\begin{document}

\begin{CJK*}{UTF8}{bkai}   %%% ZZZ %%%  <<< 在這裡更改預設中文字型、編碼
% CJK* 與 CJKspace 套件配合，會把原稿裏手動鍵入的換行鍵所引起的空白 (中文字與中文字之間) 拿掉
% 編碼：UTF8, Bg5, ...
% 中文字型名稱：TeXLive 安裝有一套明體字 bsmi 與 楷書 bkai; 其他字型視你的 LaTeX CJK 系統裝設情況而定

% 針對 latex + dvipdfmx 工作流程在 hyperref 套件的影響下，圖檔的辨識力退化
% 所作的權宜措施。可能是因為 TeXLive2007 hyperref 裏的
% 客製 graphicx / dvipdfmx 的設定檔不夠新
\ifx\myworkflow\mydvipdfmxflow
        \DeclareGraphicsExtensions{.pdf,.png,.jpg,.eps}
        \DeclareGraphicsRule{.pdf}{eps}{.xbb}{}
        \DeclareGraphicsRule{.png}{eps}{.xbb}{}
        \DeclareGraphicsRule{.jpg}{eps}{.xbb}{}
\fi

% global CJK setting
\CJKindent  %%% ZZZ %%%  段首內縮兩格

% 載入中文名詞的定義：例如，Figure -->「圖」, Chapter -->「第 x 章」
%
% this file is encoded in utf-8
% v2.02 (Sep. 12, 2012)

% 下列中文名詞的定義，如果以註解方式關閉取消，
% 則會以系統原先的預設值 (英文) 替代
% 名詞 \prechaptername 預設值為 Chapter
% 名詞 \postchaptername 預設值為空字串
% 名詞 \tablename 預設值為 Table
% 名詞 \figurename 預設值為 Figure
% 名詞 \lstlistingname 預設值為 Listing
%\renewcommand\prechaptername{第 } % 出現在每一章的開頭的「第 x 章」
%\renewcommand\postchaptername{ 章}
%\renewcommand{\tablename}{表} % 在文章中 table caption 會以「表 x」表示
%\renewcommand{\figurename}{圖} % 在文章中 figure caption 會以「圖 x」表示
%\renewcommand{\lstlistingname}{列表} % 在文章中 listing caption 會以「列表 x」表示


% 下列中文名詞的定義，用於論文固定的各部分之命名 (出現於目錄與該頁標題)
\newcommand{\nameInnerCover}{Inside Cover}
\newcommand{\nameCommitteeForm}{Committee Form}
\newcommand{\nameCopyrightForm}{Copyright Form}
\newcommand{\nameCabstract}{中文摘要}
\newcommand{\nameEabstract}{Abstract}
\newcommand{\nameAckn}{Acknowledgment}
\newcommand{\nameToc}{Table of Contents}
\newcommand{\nameLot}{List of Tables}
\newcommand{\nameTof}{List of Figures}
\renewcommand{\lstlistlistingname}{程式列表目錄}
\newcommand{\nameSlist}{名詞與符號說明}
\newcommand{\nameRef}{Reference}
\newcommand{\nameVita}{Vita}


% 是否以中文數字一、二、三呈現章別，例如「第一章」，則會以「第 1 章」呈現
% %
% this file is encoded in utf-8
% v2.02 (Sep. 12, 2012)

%%% ZZZ %%% 如果不是用 LaTeX CJKnumb 套件，
%% 則必須自己製作陽春的 \CJKnumber
%% 請把下列介於 >>> 與 <<< 之間
%% 的文字行打開 (行首移除百分號)
%% 中文數字對應指令 (只能對應 1 到 10)
%% 參考自 LaTeX-CJK doc: CJK.txt, zh-Hant.cpx
%% 用法：\CJKnumber{1} 可得一

%%% 陽春 \CJKnumber 命令定義 >>>
%\makeatletter
%\@ifundefined{CJKnumber}
%  {\def\CJKnumber#1{\ifcase#1\or一\or二\or三\or四\or五\or%
% 六\or七\or八\or九\or十\fi}}{}
%\makeatother
%%% <<< 陽春 \CJKnumber 命令定義

%%%%%%%%%%%%%%%%%%%%%%%%%%%
%%% 天干對應 \CJKordinal 
\makeatletter
\@ifundefined{CJKordinal}
  {\def\CJKordinal#1{\ifcase#1\or甲\or乙\or丙\or丁\or戊\or%
 己\or庚\or辛\or壬\or癸\fi}}{}
\makeatother
%%


%%%%%%%%%%%%%%%%%%%%%%%%%%%%%%%%%%%%


% 請依需要選擇其中一種表現方式，把它所對應的指令列打開，其他沒有用到的表現方式的對應指令列請關閉。(用行首百分號)

%% 第一種目錄格式：
%%	1  簡介 ............................ 1
%%
%%      章別 (chapter counter) 「1」前後沒有其他文字，
%%
%%      內文章標題是
%%		第 1 章	簡介
%%	\tocprechaptername, \tocpostchaptername 都設成沒有內容的空字串
%%	\tocChNumberWidth 設成 1.4em (預設)
%%      底下五行指令請打開
%\renewcommand\tocprechaptername{}
%\renewcommand\tocpostchaptername{}
%\setlength{\tocChNumberWidth}{1.4em}


%% 第二種目錄格式：
%%	一、簡介 ............................ 1
%%
%%      章別 (chapter counter) 「一」前沒有文字，後有頓號，
%%
%%      內文章標題是
%%		第一章		簡介
%%	\tocprechaptername 設成沒有內容的空字串
%%	\tocpostchaptername 設成頓號
%%	\tocChNumberWidth 設成 2em
%%      底下四行指令請打開 (預設)

%\renewcommand\countermapping[1]{\CJKnumber{#1}}
%\renewcommand\tocprechaptername{}
%\renewcommand\tocpostchaptername{、}
%\setlength{\tocChNumberWidth}{2em}
%\renewcommand\prechaptername{第} % 出現在每一章的開頭的「第x章」, x前後沒有空白
%\renewcommand\postchaptername{章}


%% 第三種目錄格式：
%%	第一章、簡介 ......................... 1
%%
%%      章別 (chapter counter) 「一」前有「第」，後有「章」與頓號，
%%      內文章標題是
%%		第一章		簡介
%%	\tocprechaptername 設成「第」
%%	\tocpostchaptername 設成「章、」
%%	\tocChNumberWidth 設成 3em
%%      底下四行指令請打開
%\renewcommand\countermapping[1]{\CJKnumber{#1}}
%\renewcommand\tocprechaptername{第}
%\renewcommand\tocpostchaptername{章、}
%\setlength{\tocChNumberWidth}{3em}
%\renewcommand\prechaptername{第} % 出現在每一章的開頭的「第x章」, x前後沒有空白
%\renewcommand\postchaptername{章}



%% 可以依照需要作彈性的設定
%%
%% 章別 (數字，包括後面的字串) 的寬度 \tocChNumberWidth，
%% 會影響章名與章別之間的間隔 (太少則相疊，太多則留白)
%% 建議設成 \tocpostchaptername 內容字數加一，做為 em 的倍數，
%% 但至少也要有 1.4 倍。


%%% 以下是載入前頁、本文、後頁
% 請勿更動
% 如需針對個別章節獨立編譯
% 請在 my_chapters.tex 檔裡對個別章節的 \input 指令以行首百分號方式做開關。
%
% this file is encoded in utf-8
% v2.02 (Sep. 12, 2012)

% #### front_pages 前頁 ####
% - ncku_frontpages.tex 此檔只提供整體架構的定義，不需更動
% - 包括封面、書名頁、中文摘要、英文摘要、誌謝、目錄、表目錄、圖目錄、符號說明
% this file is encoded in utf-8

%%%% Note %%%%%
% - 無須修改本檔內容
% - 除非校方修改了封面、書名頁、中文摘要、英文摘要、誌謝、目錄、表目錄、圖目錄、符號說明等頁之格式

% make the line spacing in effect
\renewcommand{\baselinestretch}{\mybaselinestretch}
\large % it needs a font size changing command to be effective

% 此處只是預設值，不需更改此處，請更改 texFolder/name.tex內容
\newcommand\cTitle{論文題目}
\newcommand\eTitle{THESIS TITLE}
\newcommand\myCname{中文名稱}
\newcommand\myEname{英文名稱}
\newcommand\advisorCnameA{指導教授A中文名稱}
\newcommand\advisorEnameA{指導教授A英文名稱}
\newcommand\advisorCnameB{指導教授B中文名稱}
\newcommand\advisorEnameB{指導教授B英文名稱}
\newcommand\advisorCnameC{指導教授C中文名稱}
\newcommand\advisorEnameC{指導教授C英文名稱}
\newcommand\univCname{中文大學名稱}
\newcommand\univEname{英文大學名稱}
\newcommand\deptCname{中文系所名稱}
\newcommand\deptEname{英文系所名稱}
\newcommand\collEname{英文學院全稱}
\newcommand\degreeCname{中文學位名稱, Ex : 碩士}
\newcommand\degreeEname{英文學位名稱, Ex : Master}
\newcommand\cYear{中文論文產生年份}
\newcommand\cMonth{中文論文產生月份}
\newcommand\eYear{英文論文產生年份}
\newcommand\eMonth{英文論文產生月份}
\newcommand\ePlace{英文, 論文發表地點}


 % user's names; to replace those default variable definitions
%this file is encoded in utf-8

% ### 檔案說明 ### 
% - 填入你的論文題目、姓名等資料

% 論文題目 (中文)
\renewcommand\cTitle{利用大型語言模型於低精度穿戴式感測行為辨識}

% 論文題目 (英文)
\renewcommand\eTitle{Utilizing Large Language Model for Human Activity Recognition from Coarse-Grained Wearable Sensors}

% 我的姓名 (中文)
\renewcommand\myCname{蘇可維}

% 我的姓名 (英文)
\renewcommand\myEname{Su-Ko Wei}

% 指導教授A的姓名 (中文)
\renewcommand\advisorCnameA{莊坤達博士}

% 指導教授A的姓名 (英文)
\renewcommand\advisorEnameA{Dr.~Kun-Ta Chuang}

% 指導教授B的姓名 (中文)
\renewcommand\advisorCnameB{}

% 指導教授B的姓名 (英文)
\renewcommand\advisorEnameB{}

% 指導教授C的姓名 (中文)
\renewcommand\advisorCnameC{}

% 指導教授C的姓名 (英文)
\renewcommand\advisorEnameC{}

% 校名 (中文)
\renewcommand\univCname{國立成功大學}

% 校名 (英文)
\renewcommand\univEname{National Cheng Kung University, Tainan, Taiwan, R.O.C.}

% 系所全名 (中文)
\renewcommand\deptCname{資訊工程學系}
% 系所全名 (英文)
\renewcommand\deptEname{Institute of Computer Science and Information Engineering}

% 學院全名 (英文)
\renewcommand\collEname{College of Electrical Engineering and Computer Science}

% 學位名 (中文)
\renewcommand\degreeCname{碩士}

% 學位名 (英文)
\renewcommand\degreeEname{Master of Engineering}

% 口試年份 (中文、民國)
\renewcommand\cYear{一一三}

% 口試月份 (中文)
\renewcommand\cMonth{七}

% 口試年份 (阿拉伯數字、西元)
\renewcommand\eYear{2024} 

% 口試月份 (英文)
\renewcommand\eMonth{July}

% 學校所在地 (英文)
\renewcommand\ePlace{Tainan, Taiwan}

%畢業級別；用於書背列印；若無此需要可忽略
\newcommand\GraduationClass{113}

%%%%%%%%%%%%%%%%%%%%%%


% 使用 hyperref 在 pdf 簡介欄裡填入相關資料
\ifx\hypersetup\undefined
        \relax  % do nothing
\else
        \hypersetup{
        pdftitle=\cTitle,
        pdfauthor=\myCname}
\fi
        

\newcommand\itsempty{}
%%%%%%%%%%%%%%%%%%%%%%%%%%%%%%%
%       NCKU cover 封面
%%%%%%%%%%%%%%%%%%%%%%%%%%%%%%%
\begin{titlepage}
\thispagestyle{EmptyWaterMarkPage}

\begin{center}

\makebox[12cm][s]{\Huge\univCname}\\  %顯示中文校名
\vspace{0.5cm}
\makebox[12cm][s]{\Huge\deptCname}\\ %顯示中文系所名
\vspace{1.5cm}
\makebox[6cm][s]{\Huge\degreeCname 論文}\\ %顯示論文種類 (中文)
\vspace{1.5cm}


\renewcommand{\baselinestretch}{1}   %行距 1 倍
\Large\cTitle\\  % 中文題目
\vspace{1cm}
\Large\eTitle\\ %英文題目
\vspace{2cm}

\hspace{2.4em}%
\makebox[4.8em][s]{\Large 研究生}%
\makebox[1em][c]{\Large ：}%
\makebox[7.2em][l]{\Large\myCname}\\  % 顯示作者中文名

%
\hspace{2.4em}%
\makebox[4.8em][s]{\Large 指導教授}%
\makebox[1em][c]{\Large ：}%
\makebox[7.2em][l]{\Large\advisorCnameA}\\  %顯示指導教授A中文名
%

%%%%% 判斷是否有共同指導的教授 B %%%%%
\ifx \advisorCnameB  \itsempty
\relax % 沒有 B 教授，所以不佔版面，不印任何空白
\else
% 共同指導的教授 B
\hspace{2.4em}%
\makebox[4.8em][s]{}%
\makebox[1em][c]{}%
\makebox[7.2em][l]{\Large\advisorCnameB}\\%顯示指導教授B中文名
\fi
%

%%%%% 判斷是否有共同指導的教授 C %%%%%
\ifx \advisorCnameC  \itsempty
\relax % 沒有 C 教授，所以不佔版面，不印任何空白
\else
% 共同指導的教授 C
\hspace{2.4em}%
\makebox[4.8em][s]{}%
\makebox[1em][c]{}%
\makebox[7.2em][l]{\Large\advisorCnameC}\\%顯示指導教授B中文名
\fi

%
\vfill
\makebox[10cm][s]{\Large 中華民國\cYear 年\cMonth 月}%
%
\end{center}

% Resume the line spacing to the desired setting
\renewcommand{\baselinestretch}{\mybaselinestretch}   %恢復原設定
% it needs a font size changing command to be effective
% restore the font size to normal
\normalsize
\end{titlepage}
%%%%%%%%%%%%%

%%---------------------------------------------------------------------------
%% 英文封面
%%---------------------------------------------------------------------------
%    \thispagestyle{empty}%
%   \null\vskip0.5in
%    \begin{center}
%        \hyphenpenalty=10000\Large\uppercase\expandafter{\eTitle}\\
%        (DRAFT)
%    \end{center}
%    \vfill
%    \begin{center}
%        \large\rm By \\
%        {\large\rm \myEname}
%    \end{center}
%    \vfill
%    \begin{center}
%        \footnotesize SUBMITTED IN PARTIAL FULFILLMENT OF THE \\
%        REQUIREMENTS FOR THE DEGREE OF \\
%        MASTER OF SCIENCE IN \\
%        DEPARTMENT OF COMPUTER SCIENCE AND INFORMATION ENGINEERING \\
%        NATIONAL CHENG-KUNG UNIVERSITY \\
%        TAINAN, TAIWAN, R.O.C. \\
%        {\eMonth \ \eYear}
%    \end{center}
%    \vskip0.75in
%    \begin{center}
%        \rm \copyright\ Copyright by {\myEname}, \eYear
%    \end{center}
%\newpage









%% 從摘要到本文之前的部份以小寫羅馬數字印頁碼
% 但是從「書名頁」(但不印頁碼) 就開始計算
\setcounter{page}{1}
\pagenumbering{roman}
%%%%%%%%%%%%%%%%%%%%%%%%%%%%%%%
%       書名頁 
%%%%%%%%%%%%%%%%%%%%%%%%%%%%%%%
%
\newpage

% 判斷是否要浮水印？
\ifx\mywatermark\undefined 
  \thispagestyle{empty}  % 無頁碼、無 header (無浮水印)
\else
  \thispagestyle{EmptyWaterMarkPage} % 無頁碼、有浮水印
\fi

% create an entry in table of contents for 書名頁
% \phantomsection % for hyperref to register this
% \addcontentsline{toc}{chapter}{\nameInnerCover}


%%%%%%%%%%%%%%%%%%%%%%%%%%%%%%%
%       論文口試委員審定書 (計頁碼，但不印頁碼) 
%%%%%%%%%%%%%%%%%%%%%%%%%%%%%%%
%
% insert the printed standard form when the thesis is ready to bind
% 在口試完成後，再將已簽名的審定書放入以便裝訂
% create an entry in table of contents for 審定書
% 目前送出空白頁

%didi comment

%\newpage%
%{\thispagestyle{empty}%
%\phantomsection % for hyperref to register this
%\addcontentsline{toc}{chapter}{\nameCommitteeForm}%
%\mbox{}\clearpage}

%%%%%%%%%%%%%%%%%%%%%%%%%%%%%%%
%       授權書 (計頁碼，但不印頁碼) 
%%%%%%%%%%%%%%%%%%%%%%%%%%%%%%%
%
% insert the printed standard form when the thesis is ready to bind
% 在口試完成後，再將已簽名的授權書放入以便裝訂
% create an entry in table of contents for 授權書
% 目前送出空白頁

%didi comment
%\newpage%
%{\thispagestyle{empty}%
%\phantomsection % for hyperref to register this
%\addcontentsline{toc}{chapter}{\nameCopyrightForm}%
%\mbox{}\clearpage}

%%%%%%%%%%%%%%%%%%%%%%%%%%%%%%%
%       中文摘要 
%%%%%%%%%%%%%%%%%%%%%%%%%%%%%%%
\newpage
\thispagestyle{plain}
\phantomsection % for hyperref to register this
\addcontentsline{toc}{chapter}{\nameCabstract}

\begin{center}

\renewcommand{\baselinestretch}{1}   %行距 1 倍

\large\cTitle\\  %中文題目
\vspace{0.83cm}

\makebox[3em][l]{學生：}%
\makebox[4.8em][l]{\myCname}%學生中文姓名
\hfill%
%
\makebox[5em][l]{指導教授：}%
\makebox[7.2em][l]{\advisorCnameA}\\ %教授A中文姓名
%

%%%%% 判斷是否有共同指導的教授 B %%%%%
\ifx \advisorCnameB  \itsempty
\relax % 沒有 B 教授，所以不佔版面，不印任何空白
\else
% 共同指導的教授B
\makebox[3em][l]{}%
\makebox[4.8em][l]{}%%%%%
\hfill%
\makebox[5em][l]{}%
\makebox[7.2em][l]{\advisorCnameB}\\ %教授B中文姓名
\fi

%%%%% 判斷是否有共同指導的教授 C %%%%%
\ifx \advisorCnameC  \itsempty
\relax % 沒有 C 教授，所以不佔版面，不印任何空白
\else

% 共同指導的教授C
\makebox[3em][l]{}%
\makebox[4.8em][l]{}%%%%%
\hfill%
\makebox[5em][l]{}%
\makebox[7.2em][l]{\advisorCnameC}\\ %教授B中文姓名
\fi

%
\vspace{0.42cm}
%

\univCname\deptCname\\ %校名系所名
\vspace{0.83cm}
\makebox[2.5cm][s]{摘要}\\
\end{center}

% Resume the line spacing to the desired setting
\renewcommand{\baselinestretch}{\mybaselinestretch}   %恢復原設定
\normalsize
%%%%%%%%%%%%%
人類行為辨識（HAR）是穿戴式感測技術中的重要研究領域，對於健康監測、運動分析和智慧生活應用具有重要意義。傳統的穿戴式感測器往往受限於低精度數據和計算資源限制，難以實現準確的行為辨識。隨著大型語言模型（LLM）的快速發展，其強大的序列理解和模式識別能力為解決此問題提供了新的可能。

本論文提出了一種利用大型語言模型進行低精度穿戴式感測行為辨識的創新方法。我們設計了一個適應性框架，將穿戴式感測器的粗粒度時間序列數據轉換為語言模型可理解的表示形式，並利用預訓練語言模型的強大表示學習能力來提升行為辨識的準確性。

我們的方法包含三個核心組件：（1）感測數據預處理和特徵提取模組，將原始感測信號轉換為結構化序列；（2）基於大型語言模型的行為模式學習架構，利用注意力機制捕捉時序關係；（3）多模態融合策略，整合不同感測器的信息以提升辨識性能。實驗結果顯示，相比傳統方法，我們的方法在多個公開數據集上都實現了顯著的性能提升，特別是在處理低精度和噪聲數據方面表現優異。

這項研究為穿戴式設備的智能化發展提供了新的技術路徑，推動了大型語言模型在物聯網和普適計算領域的應用創新。

關鍵字：人類行為辨識；大型語言模型；穿戴式感測器；低精度數據；時間序列分析；多模態融合

%%%%%%%%%%%%%%%%%%%%%%%%%%%%%%%
%       英文摘要 
%%%%%%%%%%%%%%%%%%%%%%%%%%%%%%%
%
\newpage
\thispagestyle{plain}

\phantomsection % for hyperref to register this
\addcontentsline{toc}{chapter}{\nameEabstract}

\begin{center}
\renewcommand{\baselinestretch}{1}   %行距 1 倍

\large\eTitle  %英文題目

\vspace{0.83cm}
\makebox[2cm][l]{Student:}%
\makebox[5cm][l]{\myEname}%學生英文姓名
\hfill%
%
\makebox[2cm][l]{Advisor:}%
\makebox[5cm][l]{\advisorEnameA}\\ %教授A英文姓名
%
% 判斷是否有共同指導的教授 B
\ifx \advisorCnameB  \itsempty
\relax % 沒有 B 教授，所以不佔版面，不印任何空白
\else
%共同指導的教授B
\makebox[2cm][l]{}%
\makebox[5cm][l]{}%%%%%
\hfill%
\makebox[2cm][l]{}%
\makebox[5cm][l]{\advisorEnameB}\\ %教授B英文姓名
\fi
%
% 判斷是否有共同指導的教授 C
\ifx \advisorCnameC  \itsempty
\relax % 沒有 C 教授，所以不佔版面，不印任何空白
\else
%共同指導的教授C
\makebox[2cm][l]{}%
\makebox[5cm][l]{}%%%%%
\hfill%
\makebox[2cm][l]{}%
\makebox[5cm][l]{\advisorEnameC}\\ %教授C英文姓名
\fi
%
\vspace{0.42cm}
Submitted to \deptEname\\
\collEname\\
\univEname\\  % 英文校名
\vspace{0.83cm}
%

ABSTRACT\\
%\vspace{0.5cm}
\end{center}

% Resume the line spacing the desired setting
\renewcommand{\baselinestretch}{\mybaselinestretch}   %恢復原設定
\normalsize
%%%%%%%%%%%%%


Human Activity Recognition (HAR) is a critical research area in wearable sensing technology, with significant implications for health monitoring, sports analysis, and smart living applications. Traditional wearable sensors are often limited by low-precision data and computational constraints, making accurate activity recognition challenging. The rapid advancement of Large Language Models (LLMs) offers new possibilities with their powerful sequence understanding and pattern recognition capabilities.

This thesis proposes an innovative approach that utilizes Large Language Models for human activity recognition from coarse-grained wearable sensors. We design an adaptive framework that transforms coarse-grained time-series data from wearable sensors into representations comprehensible by language models, leveraging the powerful representation learning capabilities of pre-trained language models to enhance activity recognition accuracy.

Our method comprises three core components: (1) a sensor data preprocessing and feature extraction module that converts raw sensor signals into structured sequences; (2) a Large Language Model-based activity pattern learning architecture that utilizes attention mechanisms to capture temporal relationships; and (3) a multi-modal fusion strategy that integrates information from different sensors to improve recognition performance. Experimental results demonstrate that our approach achieves significant performance improvements over traditional methods across multiple public datasets, particularly excelling in handling low-precision and noisy data.

This research provides a new technical pathway for the intelligent development of wearable devices and advances the application innovation of Large Language Models in the Internet of Things and ubiquitous computing domains.

Keywords: Human Activity Recognition; Large Language Models; Wearable Sensors; Coarse-grained Data; Time Series Analysis; Multi-modal Fusion


%%%%%%%%%%%%%%%%%%%%%%%%%%%%%%%
%       誌謝 (Acknowledgment)
%%%%%%%%%%%%%%%%%%%%%%%%%%%%%%%
%
\newpage
\chapter*{\protect\makebox[5cm][s]{\nameAckn}}
\phantomsection % for hyperref to register this
\addcontentsline{toc}{chapter}{\nameAckn}




%%%%%%%%%%%%%%%%%%%%%%%%%%%%%%%
%       目錄 (Category)
%%%%%%%%%%%%%%%%%%%%%%%%%%%%%%%
%
\newpage
\renewcommand{\contentsname}{\protect\makebox[5cm][s]{Contents}}
\phantomsection % for hyperref to register this
\addcontentsline{toc}{chapter}{Contents}
\tableofcontents

%%%%%%%%%%%%%%%%%%%%%%%%%%%%%%%
%       表目錄 (List of Tables)
%%%%%%%%%%%%%%%%%%%%%%%%%%%%%%%
%
\newpage
\renewcommand{\listtablename}{\protect\makebox[5cm][s]{\nameLot}}
\phantomsection % for hyperref to register this
\addcontentsline{toc}{chapter}{\nameLot}
\listoftables

%%%%%%%%%%%%%%%%%%%%%%%%%%%%%%%
%       圖目錄 (List of Figures)
%%%%%%%%%%%%%%%%%%%%%%%%%%%%%%%
%
\newpage
\renewcommand{\listfigurename}{\protect\makebox[5cm][s]{\nameTof}}
\phantomsection % for hyperref to register this
\addcontentsline{toc}{chapter}{\nameTof}
\listoffigures

%%%%%%%%%%%%%%%%%%%%%%%%%%%%%%%
%       程式列表目錄 (List of Listings)
%%%%%%%%%%%%%%%%%%%%%%%%%%%%%%%
%
%\newpage
%\phantomsection % for hyperref to register this
%\addcontentsline{toc}{chapter}{\lstlistlistingname}
%\lstlistoflistings



\newpage
%% 論文本體頁碼回復為阿拉伯數字計頁，並從頭起算
\pagenumbering{arabic}
%%%%%%%%%%%%%%%%%%%%%%%%%%%%%%%%
 

%Human Activity Recognition (HAR) is a critical research area in wearable sensing technology, with significant implications for health monitoring, sports analysis, and smart living applications. Traditional wearable sensors are often limited by low-precision data and computational constraints, making accurate activity recognition challenging. The rapid advancement of Large Language Models (LLMs) offers new possibilities with their powerful sequence understanding and pattern recognition capabilities.

This thesis proposes an innovative approach that utilizes Large Language Models for human activity recognition from coarse-grained wearable sensors. We design an adaptive framework that transforms coarse-grained time-series data from wearable sensors into representations comprehensible by language models, leveraging the powerful representation learning capabilities of pre-trained language models to enhance activity recognition accuracy.

Our method comprises three core components: (1) a sensor data preprocessing and feature extraction module that converts raw sensor signals into structured sequences; (2) a Large Language Model-based activity pattern learning architecture that utilizes attention mechanisms to capture temporal relationships; and (3) a multi-modal fusion strategy that integrates information from different sensors to improve recognition performance. Experimental results demonstrate that our approach achieves significant performance improvements over traditional methods across multiple public datasets, particularly excelling in handling low-precision and noisy data.

This research provides a new technical pathway for the intelligent development of wearable devices and advances the application innovation of Large Language Models in the Internet of Things and ubiquitous computing domains.

Keywords: Human Activity Recognition; Large Language Models; Wearable Sensors; Coarse-grained Data; Time Series Analysis; Multi-modal Fusion
\chapter{Introduction}

Human Activity Recognition (HAR) using wearable sensors has emerged as a fundamental technology for numerous applications including mobile health monitoring, sports analytics, human-computer interaction, and smart environment systems \cite{Bulling2014HARSurvey, Lara2013Mobile}. The ability to automatically identify and classify daily activities from sensor data enables continuous and unobtrusive monitoring of human behavior, providing valuable insights for healthcare, fitness tracking, and lifestyle analysis \cite{Shoaib2015Survey}.

A critical and underexplored challenge in HAR is the significant drop in recognition performance when sensor data is downsampled or collected at low temporal resolution, a scenario increasingly common in consumer-grade wearables. Reducing the sampling frequency of sensors such as accelerometers and gyroscopes can dramatically extend battery life and operational time for smart devices, making low-grain data collection highly desirable for real-world, long-term deployment. However, this comes at the cost of information loss, increased noise, and reduced discriminative power, which can severely degrade the accuracy of activity recognition models.

Traditional HAR approaches predominantly rely on motion sensors such as accelerometers, gyroscopes, and magnetometers embedded in smartphones or dedicated wearable devices \cite{Chen2021WearableHAR}. While these sensors offer the advantage of being lightweight, power-efficient, and readily available in consumer devices, they often produce coarse-grained data with limited precision compared to specialized research-grade equipment. This limitation poses significant challenges for accurate activity recognition, particularly when distinguishing between subtle variations in movement patterns or activities with similar motion characteristics.

Recent advances in deep learning have revolutionized the field of HAR, with convolutional neural networks (CNNs) and recurrent neural networks (RNNs) showing remarkable performance improvements over traditional machine learning approaches \cite{Wang2019DeepHAR, Yang2015DeepConvLSTM}. Deep learning models have demonstrated their ability to automatically extract meaningful features from raw sensor data, reducing the need for manual feature engineering \cite{Hammerla2016Deep, Ronao2016Human}. However, most existing deep learning approaches for HAR are designed for high-quality sensor data and may not perform optimally when dealing with the noisy, low-resolution signals typical of consumer-grade wearable devices \cite{Ignatov2018RealTime}.

The emergence of Large Language Models (LLMs) and transformer architectures has opened new possibilities for sequence modeling and pattern recognition across diverse domains \cite{Vaswani2017Attention}. Models such as BERT \cite{Devlin2019BERT} and GPT-3 \cite{Brown2020GPT3} have demonstrated exceptional capabilities in understanding complex sequential patterns and contextual relationships. While these models were originally designed for natural language processing tasks, their underlying transformer architecture and attention mechanisms show promise for modeling temporal sequences in other domains, including sensor data analysis.

Recent work has begun to explore the application of LLMs to sensor-based HAR, introducing frameworks for aligning language models with motion sensors for activity recognition. These approaches typically involve specialized alignment stages and task-aware tuning procedures to bridge the gap between natural language processing and sensor-based activity recognition. However, these existing methods face several critical limitations when applied to real-world deployment scenarios involving downsampled, coarse-grained wearable sensors.

The primary challenges include computational requirements and memory footprint that present significant barriers for deployment on resource-constrained wearable devices, where power efficiency and real-time processing are paramount. Additionally, existing approaches have primarily focused on high-quality sensor data from controlled laboratory settings, but have not adequately addressed the unique challenges posed by low-frequency, low-precision sensor signals that are characteristic of consumer-grade wearable devices. Furthermore, the temporal resolution and sampling rates of downsampled wearable sensors may not align optimally with the sequence processing capabilities of language models, necessitating specialized data representation and tokenization strategies.

Our work addresses these challenges through a novel approach that leverages the complementary strengths of advanced time-series processing and language model capabilities. The self-attention mechanism in transformers can effectively capture long-range dependencies in temporal sensor data, potentially identifying subtle patterns that traditional approaches might miss in noisy, low-resolution signals. The pre-training paradigm used in LLMs could enable knowledge transfer from large-scale datasets to smaller, domain-specific HAR datasets, which is particularly valuable when working with limited, downsampled sensor data. Additionally, the contextual understanding capabilities of LLMs may help in disambiguating similar activities by considering broader temporal context, which is crucial when dealing with the reduced discriminative power of low-frequency sensors.

This thesis is specifically motivated by the need to enable accurate and robust HAR on downsampled (low-grain) sensor data, with the ultimate goal of enabling longer battery life and more practical, real-world deployment of smart wearable devices.

A key contribution of this thesis is a \textbf{two-stage pipeline} that first restores high-resolution temporal detail from low-frequency sensor signals using an advanced upsampling framework and then performs activity recognition with a token-efficient downstream model. By separating data enhancement from recognition, we unlock the complementary strengths of state-of-the-art time-series imputation and sensor-language modelling, achieving robust HAR even when the raw input is severely downsampled.

The main contributions of this work include:
\begin{enumerate}
\item A novel \textbf{upsampler + downstream model} framework that integrates cutting-edge time-series imputation with large-scale sensor-language modelling for coarse-grained wearable HAR.
\item An end-to-end processing pipeline and tokenisation strategy that enables efficient fine-tuning of compact models on imputed sensor sequences, facilitating on-device inference without sacrificing accuracy.

\end{enumerate}

The remainder of this thesis is organized as follows: Chapter 2 reviews related work in human activity recognition, deep learning approaches for sensor data analysis, and the application of transformer models to time series data, with particular emphasis on recent developments in LLM-based sensor data processing and the challenges of low-grain data. Chapter 3 presents our proposed methodology, including the downsampled sensor data tokenization approach and the efficient LLM-based HAR architecture. Chapter 4 describes the experimental setup and presents comprehensive evaluation results across multiple datasets and sensor precision levels. Chapter 5 discusses the implications of our findings, limitations of the current approach, and future research directions. Finally, Chapter 6 concludes the thesis and summarizes the key contributions.
\chapter{Related Work}

\section{Human Activity Recognition with Wearable Sensors}

\hspace{2em}Human Activity Recognition (HAR) using wearable sensors has emerged as a fundamental research area with applications spanning healthcare monitoring, smart homes, and fitness tracking. Early comprehensive surveys by \cite{Bulling2014HARSurvey} and \cite{Lara2013Mobile} established the foundational understanding of sensor-based activity recognition, highlighting the challenges of feature extraction from multi-modal sensor data and the importance of temporal modeling for sequential activities.

The evolution of HAR has been significantly influenced by advances in deep learning methodologies. Traditional approaches relied heavily on hand-crafted features and shallow machine learning models, but the introduction of deep neural networks revolutionized the field. \cite{Wang2019DeepHAR} provided a comprehensive survey demonstrating how deep learning techniques could automatically learn hierarchical representations from raw sensor data, eliminating the need for manual feature engineering. This paradigm shift enabled more robust and generalizable activity recognition systems.

\subsection{Deep Learning Approaches for HAR}

\hspace{2em}The application of deep learning to HAR has yielded remarkable improvements in recognition accuracy and model generalization. \cite{Yang2015DeepConvLSTM} pioneered the use of deep convolutional neural networks on multichannel time series data, demonstrating superior performance compared to traditional methods. Their approach effectively captured both spatial relationships between sensor channels and temporal dependencies within activity sequences.

Convolutional Neural Networks (CNNs) have proven particularly effective for sensor data processing. \cite{Ignatov2018RealTime} developed real-time HAR systems using CNNs on accelerometer data, achieving impressive accuracy while maintaining computational efficiency suitable for mobile deployment. Similarly, \cite{Ronao2016Human} explored smartphone-based HAR using deep learning, showcasing the potential for ubiquitous activity monitoring through everyday devices.

The integration of recurrent architectures has further enhanced temporal modeling capabilities. \cite{Hammerla2016Deep} conducted comprehensive comparisons of deep, convolutional, and recurrent models for wearable-based HAR, establishing best practices for different types of sensor configurations and activity categories. Their work highlighted the importance of model architecture selection based on the specific characteristics of the target application.

\section{Transformer and Attention Mechanisms in HAR}

\hspace{2em}The transformer architecture, originally introduced by \cite{Vaswani2017Attention}, has revolutionized sequence modeling across various domains. The self-attention mechanism's ability to capture long-range dependencies and parallel processing capabilities have made it particularly attractive for sensor data analysis. Recent works have successfully adapted transformer architectures for HAR applications, demonstrating significant improvements over traditional recurrent approaches.

\cite{Li2024P2LHAP} introduced P2LHAP, a novel Patch-to-Label Seq2Seq framework that tackles activity recognition, segmentation, and forecasting simultaneously. Their approach divides sensor data streams into patches served as input tokens, enabling unified processing of multiple HAR tasks within a single model. The patch-based tokenization strategy has proven effective for handling variable-length sensor sequences while maintaining computational efficiency.

Recent advances have also explored specialized transformer variants for sensor data. \cite{Zhang2025MoPFormer} proposed MoPFormer, a Motion-Primitive Transformer that enhances interpretability by tokenizing sensor signals into semantically meaningful motion primitives. This approach addresses the critical challenge of cross-dataset generalization by learning fundamental movement patterns that remain consistent across different data sources.

\cite{Zhao2025ActionFormer} adapted the ActionFormer architecture for sensor-based HAR, demonstrating substantial improvements in activity boundary detection and classification accuracy. Their work highlighted the importance of detecting informative channels in multi-sensor configurations, achieving significant performance gains through attention-based channel selection mechanisms.

\section{Large Language Models for Sensor Data}

\hspace{2em}The recent success of Large Language Models (LLMs) has inspired researchers to explore their application to sensor data understanding. The foundational work in this area includes BERT \cite{Devlin2019BERT} and GPT-3 \cite{Brown2020GPT3}, which demonstrated the power of large-scale pre-training for natural language understanding. These advances have motivated the development of similar approaches for sensor modalities.

\cite{li2024sensorllm} introduced SensorLLM, a pioneering work that aligns large language models with motion sensors for human activity recognition. Their two-stage framework consists of Sensor-Language Alignment for learning shared representations between sensor data and natural language, followed by Task-Aware Tuning for specific downstream applications. SensorLLM demonstrated remarkable capabilities in zero-shot activity recognition and cross-modal reasoning, establishing a new paradigm for sensor data understanding.

Building upon this foundation, \cite{Zhang2025SensorLM} developed SensorLM, a comprehensive sensor-language foundation model that processes over 59.7 million hours of data from more than 103,000 participants. Their hierarchical caption generation pipeline captures statistical, structural, and semantic information from sensor data, enabling sophisticated natural language interactions with wearable sensor systems.

\cite{Civitarese2024ADLLLM} explored the application of LLMs for Activities of Daily Living recognition in smart home environments. Their ADL-LLM system transforms raw sensor data into textual representations processed by LLMs for zero-shot recognition, demonstrating the potential for deployment in real-world assisted living scenarios.

% Newly added works on zero-shot LLM-based HAR and comprehensive surveys
\cite{Ji2024HARGPT} presented \textit{HARGPT}, a systematic study investigating whether off-the-shelf LLMs can serve as zero-shot human activity recognizers without any gradient-based fine-tuning. By feeding raw, downsampled (\SI{10}{\hertz}) IMU sequences directly into GPT-4 and other large models with carefully crafted role-play and chain-of-thought prompts, HARGPT achieved approximately 80\% accuracy on the Capture-24 and HHAR benchmarks. This evidence reveals an emergent ability of LLMs to transform noisy time–series signals into high-level activity semantics purely through prompting, albeit at the expense of large context windows and intensive prompt engineering. In contrast, our two-stage pipeline fine-tunes a compact Chronos encoder--LLaMA architecture on teacher--student enhanced, coarse-grained data, eliminating elaborate prompt design while enabling efficient on-device inference.

Complementing these empirical findings, \cite{Ferrara2024LLMSurvey} provides the first dedicated survey on LLM-driven wearable sensing. The paper synthesizes early trends, publicly available datasets, evaluation protocols, and open challenges—such as privacy, robustness, and energy efficiency—thereby situating the rapidly growing body of LLM-for-HAR research within a broader interdisciplinary landscape.

\subsection{High-Level Reasoning with LLMs}

\hspace{2em}Beyond basic activity classification, recent research has investigated LLMs' capabilities for high-level reasoning over sensor traces. \cite{Ouyang2024LLMSense} developed LLMSense, which harnesses LLMs for complex spatiotemporal reasoning tasks such as dementia diagnosis and occupancy tracking. Their framework demonstrates how LLMs can comprehend long-term sensor patterns and make sophisticated inferences that go beyond simple activity labeling.

\cite{Sun2024IoTActivity} presented an AI-based system utilizing IoT-enabled ambient sensors and LLMs for complex activity tracking in elderly care. Their approach combines edge device processing with cloud-based LLM reasoning, addressing privacy concerns while maintaining the sophisticated reasoning capabilities required for healthcare applications.

The integration of multiple sensor modalities has been explored by \cite{Mo2024IoTLM}, who introduced IoT-LM, a large multisensory language model for the Internet of Things. Their MultiIoT dataset encompasses over 1.15 million samples from 12 modalities, enabling comprehensive multisensory understanding and interactive question-answering capabilities.

\section{Energy Efficiency and Resource Constraints}

\hspace{2em}A critical challenge in HAR deployment is the balance between recognition accuracy and computational efficiency, particularly for battery-powered wearable devices. This challenge has motivated extensive research into lightweight model architectures and energy-efficient processing strategies.

Traditional approaches to address computational constraints include model compression techniques, quantization, and pruning. However, these methods often require careful tuning and may significantly impact model performance. The development of inherently efficient architectures has emerged as a more promising direction.

\cite{Ruan2024PAT} introduced the Pretrained Actigraphy Transformer (PAT), designed specifically for time-series wearable movement data with emphasis on computational efficiency. Their lightweight transformer architecture achieves state-of-the-art performance in mental health prediction tasks while maintaining interpretability and deployment feasibility on resource-constrained devices.

\subsection{Data Sampling and Downsampling Strategies}

\hspace{2em}The relationship between sensor data resolution and recognition performance presents a fundamental trade-off in wearable HAR systems. Higher sampling rates provide richer temporal information but significantly increase computational and energy costs. Conversely, lower sampling rates extend battery life but may compromise recognition accuracy.

This trade-off becomes particularly critical for continuous monitoring applications where devices must operate for extended periods without charging. The challenge of maintaining recognition performance with downsampled sensor data remains largely unexplored in the literature, representing a significant gap that our work addresses.

Current HAR systems typically assume high-frequency sensor data (50-100 Hz or higher), which may not be practical for real-world deployment scenarios where energy efficiency is paramount. The development of effective methods for coarse-grained sensor data processing is essential for enabling practical wearable HAR systems.

\section{Research Gaps and Motivation}

\hspace{2em}Despite significant advances in HAR research, several critical gaps remain. First, while transformer-based approaches have shown promise, their application to coarse-grained sensor data has received limited attention. Most existing work assumes high-resolution sensor inputs, which may not be feasible in energy-constrained environments.

Second, the integration of LLMs with sensor data has primarily focused on high-frequency, well-sampled data. The effectiveness of LLM-based approaches for downsampled sensor data remains unexplored, despite its practical importance for real-world deployment.

Third, current research lacks comprehensive analysis of the attention patterns and tokenization strategies specifically designed for low-precision sensor data. Understanding how transformer models process coarse-grained sensor information is crucial for developing effective architectures.

Finally, there is insufficient investigation into the fundamental limits of activity recognition with reduced sensor resolution. Establishing these limits and developing methods to approach them is essential for advancing the field toward practical, energy-efficient HAR systems.

Our work addresses these gaps by developing novel tokenization strategies for downsampled sensor data, investigating LLM architectures optimized for resource-constrained environments, and providing comprehensive experimental evaluation across various downsampling levels. Through this research, we aim to bridge the gap between theoretical advances in HAR and their practical deployment in real-world wearable systems.

\input{texFolder/preliminaries}
\chapter{Methodology}

\begin{figure}[H]
    \centering
    \includegraphics[width=1\linewidth]{figs/CH3_METHOD/Method.png}
    \caption{Overview of the proposed two-stage methodology for coarse-grained sensor HAR}
    \label{fig:methodology-overview}
\end{figure}

\hspace{2em}In this chapter, we present our novel two-stage methodology for enhancing Large Language Model-based Human Activity Recognition (HAR) performance on coarse-grained wearable sensor data. Our approach addresses the fundamental challenge of maintaining recognition accuracy while operating under severe resource constraints typical of wearable devices. The methodology consists of two interconnected stages: (1) \textbf{SAITS-based Time Series Imputation} that enhances low-resolution sensor data through state-of-the-art self-attention mechanisms, and (2) \textbf{SensorLLM Training} on the enhanced data for improved HAR performance.

% =====================
% SAITS-based Imputation for Stage 1
% =====================
\section{Stage 1: SAITS-based Time Series Imputation}

\hspace{2em}The first stage of our methodology employs SAITS (Self-Attention-based Imputation for Time Series) to enhance coarse-grained sensor data. SAITS is a state-of-the-art deep learning model designed specifically for imputing missing values and enhancing the quality of multivariate time series data. Its self-attention architecture enables efficient modeling of long-range dependencies and complex temporal patterns, outperforming traditional RNN-based and Transformer-based imputation methods in both accuracy and computational efficiency.

\subsection{Motivation and Advantages}

\hspace{2em}The motivation for adopting SAITS is its superior ability to recover fine-grained temporal information from coarse or partially observed sensor data. Unlike recursive models, SAITS leverages pure self-attention mechanisms, which provide faster training, lower memory usage, and improved imputation accuracy. Empirical studies have shown that SAITS achieves 12\%--38\% lower mean absolute error (MAE) than BRITS and 7\%--39\% lower mean squared error (MSE) than NRTSI, while requiring only 15\%--30\% of the parameters of a standard Transformer.

\subsection{SAITS Architecture}

\hspace{2em}SAITS consists of multiple stacked self-attention layers, each designed to capture both local and global temporal dependencies in the input sequence. The model operates as follows:

\begin{itemize}
    \item \textbf{Input}: Multivariate time series with missing values (e.g., low-resolution or partially observed sensor data)
    \item \textbf{Masking}: A binary mask indicates observed and missing values
    \item \textbf{Self-Attention Blocks}: Each block applies multi-head self-attention to model dependencies across all time steps and features
    \item \textbf{Imputation}: The model predicts missing values by aggregating information from observed entries using learned attention weights
    \item \textbf{Loss Function}: The imputation loss is computed only on the originally missing entries, typically using mean absolute error (MAE) or mean squared error (MSE)
\end{itemize}

\subsection{Integration into the HAR Pipeline}

\hspace{2em}In our pipeline, SAITS is trained to impute high-resolution sensor data from coarse-grained or downsampled input. The enhanced (imputed) data produced by SAITS is then used as input for downstream human activity recognition (HAR) with SensorLLM. The integration process is as follows:

\begin{enumerate}
    \item \textbf{Data Preparation}: Low-resolution and high-resolution sensor data pairs are converted into the HDF5 format expected by SAITS, with appropriate masking of missing values.
    \item \textbf{SAITS Training}: The model is trained on the training split, using the low-resolution data as input and the high-resolution data as ground truth for imputation.
    \item \textbf{Imputation}: After training, SAITS imputes missing values in the low-resolution data, producing enhanced sequences with restored temporal detail.
    \item \textbf{Downstream HAR}: The imputed data is segmented, tokenized, and used to train SensorLLM for activity recognition, as described in Stage 2.
\end{enumerate}

\subsection{SAITS Implementation Details}

\hspace{2em}We use the official SAITS implementation, with data conversion handled by a custom script that prepares the HDF5 datasets. The model is run using the provided training script, and the best checkpoint (based on validation loss) is used for imputation on the test set. The imputed data is then passed to the HAR stage for further processing.

% --- Mathematical Formulation of SAITS ---
\subsection{Mathematical Formulation}

\hspace{2em}Let $\mathbf{X}\in\mathbb{R}^{T\times d}$ denote a multivariate time series of length $T$ with $d$ sensor channels and let $\mathbf{M}\in\{0,1\}^{T\times d}$ be the corresponding binary observation mask ($M_{t,f}=1$ if the value $X_{t,f}$ is observed). SAITS processes $[\mathbf{X},\mathbf{M}]$ through three consecutive \emph{dual-masked self-attention} (DMSA) blocks.

\begin{enumerate}
    \item \textbf{Block~1 (coarse estimation)} projects the concatenated input into the model space via a linear map $\phi$, adds positional encodings, and passes the result through $G$ groups of $L$ stacked multi-head self-attention layers with shared parameters:
    \[
        \mathbf{H}^{(1)}=\operatorname{DMSA}_1\!\bigl(\phi([\mathbf{X},\mathbf{M}])\bigr),\qquad
        \mathbf{\tilde X}^{(1)}=\psi\bigl(\mathbf{H}^{(1)}\bigr),
    \]
    where $\psi$ maps the hidden representation back to the original feature dimension.  A first imputation is obtained as $\mathbf{X}'=\mathbf{M}\odot\mathbf{X}+(1-\mathbf{M})\odot\mathbf{\tilde X}^{(1)}$.
    \item \textbf{Block~2 (refinement)} applies another DMSA stack to $[\mathbf{X}',\mathbf{M}]$ to produce a refined estimate $\mathbf{\tilde X}^{(2)}$.
    \item \textbf{Block~3 (gated fusion)} fuses the two estimates via a learnable gate $\boldsymbol{\eta}\in[0,1]^{T\times d}$ that depends on the missing mask and the attention map $\mathbf{A}$ from Block~2:
    \[
        \boldsymbol{\eta}=\sigma\bigl(W[\mathbf{M},\mathbf{A}]\bigr),\qquad
        \mathbf{\tilde X}^{(3)}=(1-\boldsymbol{\eta})\odot\mathbf{\tilde X}^{(2)}+\boldsymbol{\eta}\odot\mathbf{\tilde X}^{(1)}.
    \]
\end{enumerate}

\noindent The final imputation is $\hat{\mathbf{X}}=\mathbf{M}\odot\mathbf{X}+(1-\mathbf{M})\odot\mathbf{\tilde X}^{(3)}$, which is forwarded to downstream HAR.

\paragraph{Training Objective.}  SAITS jointly minimises a reconstruction loss on observed values (Observed–Reconstruction Task, ORT) and an imputation loss on artificially masked values (Masked–Imputation Task, MIT):
\[
\mathcal{L}_{\text{SAITS}}=
\lambda_{\text{rec}}\,\operatorname{MAE}\bigl(\mathbf{M}\odot\mathbf{\tilde X}^{(3)},\,\mathbf{M}\odot\mathbf{X}\bigr)+
\lambda_{\text{imp}}\,\operatorname{MAE}\bigl(\mathbf{M}^{\mathrm h}\odot\mathbf{\tilde X}^{(3)},\,\mathbf{M}^{\mathrm h}\odot\mathbf{X}\bigr),
\]
where $\mathbf{M}^{\mathrm h}$ is a random hold-out mask used only during training and $(\lambda_{\text{rec}},\lambda_{\text{imp}})=(1,1)$ by default.  Following the original paper, the same loss is also computed on $\mathbf{\tilde X}^{(1)}$ and $\mathbf{\tilde X}^{(2)}$ and then averaged.

\paragraph{Hyper-parameter Choices.}  Unless stated otherwise we adopt the ``best'' configuration recommended by the authors: $G{=}5$, $L{=}1$, hidden dimension $d_{\text{model}}{=}256$, feed-forward dimension $d_{\text{inner}}{=}512$, $n_{\text{heads}}{=}8$, key/value dimension $d_k{=}d_v{=}32$, dropout~$0$, and the \textit{inner\_group} parameter-sharing strategy.  This yields a compact model with about $3.2\,$M parameters.

\paragraph{Positioning and Novelty.}  Prior works \cite{Du2023SAITS, li2024sensorllm} examine imputation and sensor-language modelling in isolation. In contrast, we explicitly \emph{chain} SAITS with SensorLLM and show that high-quality reconstruction is a \emph{prerequisite} for reliable LLM-based HAR under aggressive downsampling. This simple yet unexplored coupling turns out to be surprisingly effective and constitutes the core novelty of our framework.

% Algorithm: Stage 1 SAITS-based Time Series Imputation
\begin{algorithm}[t]
\caption{Stage 1 SAITS-based Time Series Imputation}
\KwIn{Low-resolution dataset $\mathcal{D}_{low}$, high-resolution dataset $\mathcal{D}_{high}$, number of DMSA groups $G$, layers per group $L$, model dimension $d_{model}$, learning rate $\eta$, training epochs $N_{ep}$}
\KwOut{Trained SAITS model $\Theta_{SAITS}$ for data enhancement}

Initialize SAITS model with $G$ DMSA groups and dimension $d_{model}$\;
Initialize Adam optimizer with learning rate $\eta$\;
Compute z-score normalization parameters $\mu$, $\sigma$ from training data\;

\For{$epoch \leftarrow 1$ \KwTo $N_{ep}$}{
    \ForEach{mini-batch $(\mathbf{X}_{low}, \mathbf{X}_{high}) \sim (\mathcal{D}_{low}, \mathcal{D}_{high})$}{
        $\mathbf{M} \leftarrow$ \textbf{GenerateMask}$(\mathbf{X}_{low}, \mathbf{X}_{high})$\;
        $\mathbf{M}^h \leftarrow$ \textbf{GenerateHoldoutMask}$(\mathbf{M})$\;
        
        % Block 1: Coarse Estimation
        $\mathbf{H}^{(1)} \leftarrow \text{DMSA}_1(\phi([\mathbf{X}_{low}, \mathbf{M}]))$\;
        $\mathbf{\tilde{X}}^{(1)} \leftarrow \psi(\mathbf{H}^{(1)})$\;
        $\mathbf{X}' \leftarrow \mathbf{M} \odot \mathbf{X}_{low} + (1-\mathbf{M}) \odot \mathbf{\tilde{X}}^{(1)}$\;
        
        % Block 2: Refinement
        $\mathbf{\tilde{X}}^{(2)} \leftarrow \text{DMSA}_2([\mathbf{X}', \mathbf{M}])$\;
        
        % Block 3: Gated Fusion
        $\mathbf{A} \leftarrow$ \textbf{GetAttentionMap}$(\text{DMSA}_2)$\;
        $\boldsymbol{\eta} \leftarrow \sigma(W[\mathbf{M}, \mathbf{A}])$\;
        $\mathbf{\tilde{X}}^{(3)} \leftarrow (1-\boldsymbol{\eta}) \odot \mathbf{\tilde{X}}^{(2)} + \boldsymbol{\eta} \odot \mathbf{\tilde{X}}^{(1)}$\;
        
        % Compute Loss
        $\mathcal{L}_{rec} \leftarrow \text{MAE}(\mathbf{M} \odot \mathbf{\tilde{X}}^{(3)}, \mathbf{M} \odot \mathbf{X}_{high})$\;
        $\mathcal{L}_{imp} \leftarrow \text{MAE}(\mathbf{M}^h \odot \mathbf{\tilde{X}}^{(3)}, \mathbf{M}^h \odot \mathbf{X}_{high})$\;
        $\mathcal{L}_{total} \leftarrow \lambda_{rec}\mathcal{L}_{rec} + \lambda_{imp}\mathcal{L}_{imp}$\;
        
        Backpropagate and update parameters to minimise $\mathcal{L}_{total}$\;
    }
    Evaluate validation MAE; apply early stopping if no improvement for 30 epochs\;
}

\Return{$\Theta_{SAITS}$}\;
\end{algorithm}

% --- Integration of Stage 1 and Stage 2 ---
\paragraph{Stage 1 to Stage 2 Integration.}
\hspace{2em}The output of Stage 1, consisting of enhanced (imputed) high-resolution sensor data, serves as the direct input to Stage 2. Specifically, the SAITS model takes low-resolution or downsampled sensor data and reconstructs temporally rich, high-resolution sequences. These enhanced sequences are then formatted and segmented as required for the SensorLLM pipeline. This integration ensures that the downstream HAR model receives input data with restored temporal detail, which is critical for accurate activity recognition, especially under aggressive downsampling or missing data scenarios. The imputed data is thus not only a pre-processing step but a crucial enabler for robust LLM-based HAR.

\section{Stage 2: SensorLLM Training for Human Activity Recognition}

\hspace{2em}The second stage of our methodology leverages SensorLLM, a framework designed to align large language models with time-series sensor data for Human Activity Recognition. This stage takes the enhanced, high-resolution data imputed by SAITS and uses it to directly train a powerful, context-aware HAR classifier.

\subsection{SensorLLM Framework}
\hspace{2em}The SensorLLM architecture is built upon three core components:
\begin{itemize}
    \item \textbf{A Pretrained Time-Series (TS) Embedder:} We use a frozen, pretrained Chronos model ($\Phi_{TS}$) as the time-series encoder. It is responsible for extracting rich temporal features from the input sensor data, converting raw time-series signals $\mathbf{X} \in \mathbb{R}^{T \times d}$ into meaningful embeddings $\mathbf{E}_{sensor} \in \mathbb{R}^{L \times d_{TS}}$.
    \item \textbf{A Pretrained Large Language Model (LLM):} A frozen, large-scale language model, LLaMA-3 ($\Phi_{LLM}$), serves as the reasoning backbone of the framework, processing both textual information and the sensor embeddings.
    \item \textbf{An Alignment Module:} A lightweight, trainable Multi-Layer Perceptron (MLP), denoted $\Theta_{Align}$, acts as a bridge. It projects the sensor embeddings from the TS embedder's space into the LLM's representation space: $\mathbf{E}_{aligned} = \Theta_{Align}(\mathbf{E}_{sensor})$, where $\mathbf{E}_{aligned} \in \mathbb{R}^{L \times d_{LLM}}$. This is a primary component updated during training.
\end{itemize}
By keeping the large TS embedder and LLM frozen, SensorLLM achieves high computational efficiency, with only a small fraction of parameters being trainable.

\subsection{Task-Aware HAR Training and Objective}
\hspace{2em}The SensorLLM model is trained directly for the HAR task using the enhanced data from Stage 1.

\paragraph{Input Representation.} For a given sensor sample $\mathbf{X}_i$ and a text prompt $Q_i$ (e.g., ``What activity is this?''), the LLM input $\mathbf{E}_{\text{input}}$ is constructed by concatenating the embeddings of the text prompt and the projected sensor data:
\[
\mathbf{E}_{\text{input}, i} = [\mathbf{E}_{\text{prompt}, i}, \mathbf{E}_{\text{aligned}, i}] = [\text{Embed}(Q_i), \Theta_{\text{Align}}(\Phi_{\text{TS}}(\mathbf{X}_i))],
\]
where $\text{Embed}(\cdot)$ is the LLM's native text embedding layer.

\paragraph{HAR Prediction.} The combined embedding is processed by the LLM to produce a final hidden state, $\mathbf{h}_{\text{last}, i} = \Phi_{\text{LLM}}(\mathbf{E}_{\text{input}, i})$. A trainable linear classifier ($\Theta_{\text{Clf}}$) then maps this state to a probability distribution over the $K$ activity classes:
\[
\hat{\mathbf{y}}_i = \text{Softmax}(\Theta_{\text{Clf}}(\mathbf{h}_{\text{last}, i})).
\]

\paragraph{Training Objective.} To counteract class imbalance, we employ a weighted cross-entropy loss. Let $\mathbf{y}_i$ be the one-hot encoded true label for sample $i$, and $w_k$ be the weight for class $c$. The loss for a single sample is:
\[
\mathcal{L}_{\text{HAR}}(\hat{\mathbf{y}}_i, \mathbf{y}_i) = - \sum_{k=1}^{K} w_k y_{ik} \log(\hat{y}_{ik}).
\]
The total loss over a mini-batch of size $B$ is $\mathcal{L} = \frac{1}{B}\sum_{i=1}^{B} \mathcal{L}_{\text{HAR}}(\hat{\mathbf{y}}_i, \mathbf{y}_i)$. The weights are typically computed as the inverse of class frequencies. During this phase, the trainable parameters are $\Theta_{\text{Align}}$ and $\Theta_{\text{Clf}}$.

\begin{algorithm}[t]
\caption{Stage 2 SensorLLM Training for HAR}
\KwIn{Enhanced dataset $\mathcal{D}_{enhanced}$ from Stage 1, Pretrained TS Embedder $\Phi_{TS}$, Pretrained LLM $\Phi_{LLM}$}
\KwOut{Trained SensorLLM for HAR: Alignment Module $\Theta_{Align}$ and Classifier $\Theta_{Clf}$}

Initialize Alignment Module $\Theta_{Align}$ and Classifier $\Theta_{Clf}$\;
\For{each epoch}{
    \ForEach{mini-batch $(\mathbf{X}, y_{true}) \sim \mathcal{D}_{enhanced}$}{
        $\mathbf{E}_{sensor} \leftarrow \Phi_{TS}(\mathbf{X})$\;
        $\mathbf{E}_{aligned} \leftarrow \Theta_{Align}(\mathbf{E}_{sensor})$\;
        $\mathbf{E}_{input} \leftarrow \textbf{Combine}(\textbf{Embed}(\text{''What activity?''}), \mathbf{E}_{aligned})$\;
        $\mathbf{h}_{last} \leftarrow \Phi_{LLM}(\mathbf{E}_{input})$\;
        $\hat{\mathbf{y}} \leftarrow \text{Softmax}(\Theta_{Clf}(\mathbf{h}_{last}))$\;
        $\mathcal{L}_{HAR} \leftarrow \text{WeightedCrossEntropyLoss}(\hat{\mathbf{y}}, y_{true})$\;
        Backpropagate and update $\Theta_{Align}$ and $\Theta_{Clf}$\;
    }
}
\Return{$\Theta_{Align}, \Theta_{Clf}$}\;
\end{algorithm}


\section{Integration and Evaluation}

\hspace{2em}The integration of our two-stage approach creates a comprehensive pipeline that transforms coarse-grained sensor data into high-quality HAR predictions. The enhanced data from Stage 1 provides the temporal richness necessary for effective Stage 2 training, while the SensorLLM framework in Stage 2 leverages both sensor patterns and contextual understanding for robust activity recognition.

\hspace{2em}This methodology enables HAR systems to achieve high performance even when operating under severe resource constraints, making it particularly suitable for deployment on consumer wearable devices where battery life and computational efficiency are paramount concerns.

\chapter{Experiments}

\section{Datasets}

\subsection{Description}
\hspace{2em}We used two public datasets in our experiments. The first dataset, SleepAccel, was collected using an Apple Watch and includes motion data, heart rate, and sleep labels from PSG. Data were collected at the University of Michigan from June 2017 to March 2019, and there are 31 subjects in total. The second dataset, DREAMT, was collected using an Empatica E4 wristband. A total of 100 unique participants were recruited from the Duke University Health System (DUHS) Sleep Disorder Lab to participate in the study between May 2022 and September 2022. The total counts and distribution of the data are shown in the table below.

\begin{table}[h]
\centering
\caption{Label Counts and Distribution for SleepAccel and DREAMT Datasets}
\begin{tabular}{|c|c|c|c|c|c|}
\hline
\textbf{Dataset} & \textbf{WAKE} & \textbf{N1} & \textbf{N2} & \textbf{N3} & \textbf{REM} \\ \hline
\multicolumn{6}{|c|}{\textbf{Counts}} \\ \hline
SleepAccel       & 2,133         & 1,624       & 12,184       & 3,189       & 5,427           \\ \hline
DREAMT           & 19,971        & 8,819       & 39,687      & 2,661       & 8,331           \\ \hline
\multicolumn{6}{|c|}{\textbf{Percentage (\%)}} \\ \hline
SleepAccel       & 8.7\%       & 6.6\%     & 49.6\%     & 13\%     & 22.1\%          \\ \hline
DREAMT           & 25.1\%       & 11.1\%     & 49.9\%     & 3.3\%     & 10.5\%          \\ \hline
\end{tabular}
\end{table}

\subsection{Preprocessing}

\hspace{2em}Due to the different sampling frequencies of sensors in the datasets, such as the heart rate at 64 Hz and motion at 32 Hz in the DREAMT dataset, and heart rate at 0.2 Hz and motion at 50 Hz in the SleepAccel dataset, it was necessary to align these data for analysis. To synchronize the data and preserve more temporal information, interpolation was used to upsample all data to 64 Hz. However, this process can introduce some artifacts, so future work may consider a compromise approach, such as upsampling some data while downsampling others.

Additionally, because the data includes preparation and measurement start stages, it is possible to accurately determine the start time of measurement for users. Therefore, in addition to three-axis accelerometer data and heart rate data, a relative time difference to the start of measurement (or lights-off) was added as a feature.

To further preprocess the data, Interquartile Range (IQR) normalization was applied to the 3-axis acceleration data and heart rate data. IQR normalization helps to reduce the impact of outliers by scaling the data based on the interquartile range, thus ensuring a more robust and stable input for the model.

\section{Baselines}

\hspace{2em}We compare our proposed model with one baseline method that aim to reduce computational costs, more baselines would be further implemented:

\begin{itemize}
    \item \textbf{CNN-Transformer}: This model uses four convolutional layers followed by the encoder component of a transformer. It is designed to leverage both convolutional and transformer architectures to achieve efficient feature extraction and sequence modeling.
\end{itemize}


\section{Experiment Setting}

\hspace{2em}We set the hyperparameter \(\beta\) in the Class-Balanced Softmax Cross Entropy to 0.99 and the confidence threshold \(\theta\) to 0.7. The Adam optimizer is used with a learning rate of 0.01, and the batch size is set to 512.

To ensure robust model validation and to prevent overfitting, we employ stratified 5-fold cross-validation. This technique involves splitting the training data into five equal parts, or folds, while maintaining the class distribution within each fold. The model is trained on four folds and validated on the remaining fold. This process is repeated five times, with each fold used exactly once for validation. The cross-validation process helps in assessing the model’s performance more reliably, as it ensures that each data point has been used for both training and validation, thereby providing a comprehensive evaluation of the model’s ability to generalize to unseen data.

\begin{table}[h]
\centering
\caption{Experiment Settings}
\begin{tabular}{|l|c|}
\hline
\textbf{Setting} & \textbf{Value} \\ \hline
\(\beta\) in Class-Balanced Softmax & 0.99 \\ \hline
Cross-Validation & Stratified 5-Fold \\ \hline
Learning Rate & 0.01 \\ \hline
Batch Size & 512 \\ \hline
Confidence Threshold \(\theta\) & 0.8 \\ \hline
\end{tabular}
\end{table}

\section{Evaluation Metrics}

\hspace{2em}To comprehensively evaluate the performance of the proposed model, we employ several evaluation metrics that capture different aspects of model effectiveness and efficiency. These metrics include accuracy, macro F1-score, Cohen's kappa, FLOPs (Floating Point Operations), and the number of model parameters.

\subsection{Accuracy}
Accuracy measures the proportion of correctly classified instances among the total instances. It is a straightforward metric but can be misleading in the case of imbalanced datasets. Accuracy is defined as:

\[
\text{Accuracy} = \frac{TP + TN}{TP + TN + FP + FN}
\]

where \(TP\) stands for True Positives, \(TN\) for True Negatives, \(FP\) for False Positives, and \(FN\) for False Negatives.

\subsection{Macro F1-Score}
\hspace{2em}The macro F1-score is the harmonic mean of precision and recall, calculated for each class independently and then averaged. This metric gives equal weight to each class, making it particularly useful for evaluating performance on imbalanced datasets. The macro F1-score is defined as:

\[
\text{F1-score} = \frac{2 \cdot \text{Precision} \cdot \text{Recall}}{\text{Precision} + \text{Recall}}
\]

\[
\text{Macro F1-score} = \frac{1}{C} \sum_{i=1}^{C} \text{F1-score}_i
\]

where \(C\) is the number of classes.

\subsection{Cohen's Kappa}
\hspace{2em}Cohen's kappa measures the agreement between predicted and true labels, taking into account the possibility of agreement occurring by chance. It is particularly useful for imbalanced datasets. Cohen's kappa is defined as:

\[
\kappa = \frac{p_o - p_e}{1 - p_e}
\]

where \(p_o\) is the observed agreement, and \(p_e\) is the expected agreement by chance.

\subsection{FLOPs (Floating Point Operations)}

\hspace{2em}FLOPs measure the computational complexity of the model by counting the number of floating-point operations required to make a prediction. This metric is crucial for evaluating the computational efficiency and feasibility of deploying the model on resource-constrained devices.

\subsection{Model Parameters}
\hspace{2em}The number of model parameters is an important metric for understanding the model's complexity and potential overfitting. A model with fewer parameters is generally more efficient and suitable for deployment on devices with limited computational resources.

By using these evaluation metrics, we aim to provide a comprehensive assessment of the model's performance, balancing accuracy, robustness, and computational efficiency to ensure its practical applicability for sleep stage classification on consumer wearable devices.

\section{Performance Comparison}

\subsection{Performance Result}

\begin{table}[ht]
    \centering
    \begin{tabular}{lcccccc}
        \toprule
        Method (\(\theta\)) & Accuracy & F1 Score & Cohen's Kappa & FLOPs (M) & Parameters (K) \\
        \midrule
        CNNTransformer & \textbf{0.787} & \textbf{0.7133} & \textbf{0.6932} & 9.8566 & 384.837  \\
        Proposed Method (0.8) & 0.7267 & 0.6476 & 0.6073 & \textbf{2.0104} & \textbf{32.01} \\
        \bottomrule
    \end{tabular}
    \caption{Comparison of performance metrics between baseline and proposed methods on the SleepAccel dataset.}
    \label{tab:performance_comparison_sleepaccel}
\end{table}

\begin{table}[ht]
    \centering
    \begin{tabular}{lcccccc}
        \toprule
        Method (\(\theta\)) & Accuracy & F1 Score & Cohen's Kappa & FLOPs (M) & Parameters (K) \\
        \midrule
        CNNTransformer & \textbf{0.697} & \textbf{0.6178} & \textbf{0.54} & 9.8566 & 384.837  \\
        Proposed Method (0.8) & 0.6308 & 0.5377 & 0.4479 & \textbf{2.0104} & \textbf{32.01} \\
        \bottomrule
    \end{tabular}
    \caption{Comparison of performance metrics between baseline and proposed methods on the DREAMT dataset.}
    \label{tab:performance_comparison_dreamt}
\end{table}

% Tables \ref{tab:performance_comparison_sleepaccel} and \ref{tab:performance_comparison_dreamt} compare the performance of the CNNTransformer baseline model with the proposed method on two different datasets: SleepAccel and DREAMT. The results indicate that the proposed method successfully reduces the computational cost by approximately 80\% and the number of parameters by nearly 90\%. Despite this significant reduction in computational complexity, there is a notable trade-off in performance metrics such as accuracy, F1 score, and Cohen's kappa.

% Figure \ref{fig:performance-combine-simple-deep} illustrates the confusion matrices of the proposed method on the SleepAccel dataset. The confusion matrices are presented for the combined output, the simple classifier, and the deep classifier. The simple classifier demonstrates high accuracy, showing its effectiveness in reducing computational costs while maintaining reasonable performance. However, the performance for the N1 sleep stage is particularly poor due to the limited number of samples available for this stage, which affects the model's ability to generalize effectively. The deep classifier is employed to capture more intricate features, improving overall performance in more complex cases.

\begin{figure}
    \centering
    \includegraphics[width=1\linewidth]{figs/performance-combine-simple-deep.png}
    \caption{Confusion matrix of the proposed method on the SleepAccel dataset: (a) Combined output (b) Simple Classifier output (c) Deep Classifier output}
    \label{fig:performance-combine-simple-deep}
\end{figure}

\subsection{Performance on Different Confidence Thresholds}

\begin{table}[ht]
    \centering
    \begin{tabular}{lcccccc}
        \toprule
        Method (\(\theta\)) & Accuracy & F1 Score & Cohen's Kappa & FLOPs (M) & Parameters (K) \\
        \midrule
        COPS-CD-1 & \textbf{0.7423} & \textbf{0.6643} & \textbf{0.6278} & 2.0330 & \textbf{31.685}\\
        COPS-CD-0.9 & 0.7354 & 0.6557 & 0.6173 & 2.0256 & 32.01 \\
        COPS-CD-0.8 & 0.7391 & 0.6576 & 0.6227 & 2.0098 & 32.01 \\
        COPS-CD-0.7 & 0.7276 & 0.6496 & 0.6094 & 1.9705 & 32.01 \\
        COPS-CD-0.6 & 0.7130 & 0.6319 & 0.5890 & \textbf{1.8962} & 32.01 \\
        \bottomrule
    \end{tabular}
    \caption{Comparison of performance metrics at different confidence thresholds on the SleepAccel dataset.}
    \label{tab:performance_of_different_confidence}
\end{table}

\begin{table}[ht]
    \centering
    \begin{tabular}{lcccccc}
        \toprule
        Method (\(\theta\)) & Accuracy & F1 Score & Cohen's Kappa & FLOPs (M) & Parameters (K) \\
        \midrule
        COPS-CD-1 & \textbf{0.6273} & \textbf{0.5275} & \textbf{0.4353} & 2.0330 & \textbf{31.685}\\
        COPS-CD-0.9 & 0.6227 & 0.5187 & 0.4325 & 2.0263 & 32.01 \\
        COPS-CD-0.8 & 0.6197 & 0.5166 & 0.4302 & 2.0157 & 32.01 \\
        COPS-CD-0.7 & 0.6223 & 0.5246 & \textbf{0.4353} & 2.0034 & 32.01 \\
        COPS-CD-0.6 & 0.6201 & 0.5189 & 0.4326 & \textbf{1.9812} & 32.01 \\
        \bottomrule
    \end{tabular}
    \caption{Comparison of performance metrics at different confidence thresholds on the DREAMT dataset.}
    \label{tab:performance_of_different_confidence}
\end{table}

\begin{table}[ht]
    \centering
    \begin{tabular}{lcccccc}
        \toprule
        Method (\(\theta\)) & Accuracy & F1 Score & Cohen's Kappa & FLOPs (M) & Parameters (K) \\
        \midrule
        CDL-CD-1 & 0.7884 & \textbf{0.7138} & 0.6948 & 9.8566 & \textbf{384.837}\\
        CDL-CD-0.9 & \textbf{0.7907} & 0.7123 & \textbf{0.7019} & 9.1792 & 386.452 \\
        CDL-CD-0.8 & 0.7687 & 0.6925 & 0.6727 & 8.8275 & 386.452 \\
        CDL-CD-0.7 & 0.7562 & 0.6734 & 0.6538 & 8.3180 & 386.452 \\
        CDL-CD-0.6 & 0.7077 & 0.6167 & 0.5866 & \textbf{7.6957} & 386.452 \\
        \bottomrule
    \end{tabular}
    \caption{Comparison of performance metrics at different confidence thresholds for the CNNTransroemer on the SleepAccel dataset.}
    \label{tab:performance_of_different_confidence}
\end{table}

\begin{table}[ht]
    \centering
    \begin{tabular}{lcccccc}
        \toprule
        Method (\(\theta\)) & Accuracy & F1 Score & Cohen's Kappa & FLOPs (M) & Parameters (K) \\
        \midrule
        CDL-CD-1 & 0.6943 & 0.6202 & 0.5429 & 9.8566 & 384.837\\
        CDL-CD-0.9 & 0.6812 & 0.6006 & 0.5370 & 9.4794 & 386.452 \\
        CDL-CD-0.8 & 0.6728 & 0.5801 & 0.5246 & 9.2373 & 386.452 \\
        CDL-CD-0.7 & 0.6337 & 0.5443 & 0.4789 & 8.9522 & 386.452 \\
        CDL-CD-0.6 & 0.6248 & 0.5370 & 0.4652 & 8.5095 & 386.452 \\
        \bottomrule
    \end{tabular}
    \caption{Comparison of performance metrics at different confidence thresholds for the CNNTransroemer on the DREAMT dataset.}
    \label{tab:performance_of_different_confidence}
\end{table}


% Tables \ref{tab:performance_of_different_confidence} 


\subsection{Implementation of Confidence Threshold on Baseline}

% \hspace{2em} Table \ref{tab:performance_of_baseline_with_confidence} 

\begin{table}[ht]
    \centering
    \begin{tabular}{lcccccc}
        \toprule
        Method (\(\theta\)) & Accuracy & F1 Score & Cohen's Kappa & FLOPs (M) & Parameters (K) \\
        \midrule
        CNNTransformer & 0.7870 & 0.7133 & 0.6932 & 9.8566 & \textbf{384.837}\\
        CNNTransformer (0.8) & \textbf{0.7978} & \textbf{0.7235} & \textbf{0.7071} & \textbf{8.9325} & 386.452 \\
        \bottomrule
    \end{tabular}
    \caption{Comparison of performance metrics at different confidence thresholds for the CNNTransformer baseline.}
    \label{tab:performance_of_baseline_with_confidence}
\end{table}



\begin{table}[ht]
    \centering
    \begin{tabular}{lcccccc}
        \toprule
        Method (\(\theta\)) & Accuracy & F1 Score & Cohen's Kappa & FLOPs (M) & Parameters (K) \\
        \midrule
        CDL-Deeponly & 0.7944 & 0.7220 & 0.7030 & \textbf{9.8566} & \textbf{384.837}\\
        CDL-CD-2 & 0.7963 & 0.7264 & 0.7072 & \textbf{9.8566} & \textbf{384.837}\\
        CDL-CD-1 & \textbf{0.7974} & \textbf{0.7282} & \textbf{0.7084} & \textbf{9.8566} & \textbf{384.837}\\
        CDL-CD-0.9 & 0.7919 & 0.7211 & 0.7011 & 9.8570 & 385.162 \\
        CDL-CD-0.8 & 0.7849 & 0.7126 & 0.6925 & 9.8570 & 385.162 \\
        CDL-CD-0.7 & 0.7801 & 0.7072 & 0.6853 & 9.8570 & 385.162 \\
        CDL-CD-0.6 & 0.7671 & 0.6921 & 0.6649 & 9.8570 & 385.162 \\
        \bottomrule
    \end{tabular}
    \caption{Comparison of performance metrics at different confidence thresholds for the CNNTransroemer-CD-one on the SleepAccel dataset.}
    \label{tab:performance_of_different_confidence}
\end{table}

\begin{table}[ht]
    \centering
    \begin{tabular}{lcccccc}
        \toprule
        Method (\(\theta\)) & Accuracy & F1 Score & Cohen's Kappa & FLOPs (M) & Parameters (K) \\
        \midrule
        CDL-Deeponly & 0.7944 & 0.7220 & 0.7030 & \textbf{9.8566} & \textbf{384.837}\\
        CDL-CD-2 & \textbf{0.8017} & 0.7312 & \textbf{0.7137} & \textbf{9.8566} & \textbf{384.837}\\
        CDL-CD-1 & 0.7951 & \textbf{0.7263} & 0.7049 & \textbf{9.8566} & \textbf{384.837}\\
        CDL-CD-0.9 & 0.7893 & 0.7151 & 0.6991 & 9.8573 & 385.482 \\
        CDL-CD-0.8 & 0.7868 & 0.7105 & 0.6949 & 9.8573 & 385.482 \\
        CDL-CD-0.7 & 0.7799 & 0.7033 & 0.6853 & 9.8573 & 385.482 \\
        CDL-CD-0.6 & 0.7657 & 0.6861 & 0.6638 & 9.8573 & 385.482 \\
        \bottomrule
    \end{tabular}
    \caption{Comparison of performance metrics at different confidence thresholds for the CNNTransroemer-CD-two on the SleepAccel dataset.}
    \label{tab:performance_of_different_confidence}
\end{table}

\begin{table}[ht]
    \centering
    \begin{tabular}{lcccccc}
        \toprule
        Method (\(\theta\)) & Accuracy & F1 Score & Cohen's Kappa & FLOPs (M) & Parameters (K) \\
        \midrule
        CDL-Deeponly & 0.7944 & 0.7220 & 0.7030 & \textbf{9.8566} & \textbf{384.837}\\
        CDL-CD-2 & 0.7944 & 0.7265 & 0.7086 & \textbf{9.8566} & \textbf{384.837}\\
        CDL-CD-1 & \textbf{0.7980} & \textbf{0.7271} & \textbf{0.7091} & \textbf{9.8566} & \textbf{384.837}\\
        CDL-CD-0.9 & 0.7978 & 0.7166 & 0.6984 & 9.8573 & 385.482 \\
        CDL-CD-0.8 & 0.7889 & 0.7097 & 0.6955 & 9.8573 & 385.482 \\
        CDL-CD-0.7 & 0.7821 & 0.7083 & 0.6876 & 9.8573 & 385.482 \\
        CDL-CD-0.6 & 0.7589 & 0.6808 & 0.6547 & 9.8573 & 385.482 \\
        \bottomrule
    \end{tabular}
    \caption{Comparison of performance metrics at different confidence thresholds for the CNNTransroemer-CD-three on the SleepAccel dataset.}
    \label{tab:performance_of_different_confidence}
\end{table}

\begin{table}[ht]
    \centering
    \begin{tabular}{lcccccc}
        \toprule
        Method (\(\theta\)) & Accuracy & F1 Score & Cohen's Kappa & FLOPs (M) & Parameters (K) \\
        \midrule
        CDL-Deeponly & 0.6965 & 0.6207 & 0.5450 & \textbf{9.8566} & \textbf{384.837}\\
        CDL-CD-2 & 0.7003 & 0.6307 & 0.5555 & \textbf{9.8566} & \textbf{384.837}\\
        CDL-CD-1 & 0.7037 & 0.6343 & 0.5594 & \textbf{9.8566} & \textbf{384.837}\\
        CDL-CD-0.9 & 0.6914 & 0.6125 & 0.5428 & 9.8570 & 385.162 \\
        CDL-CD-0.8 & 0.6952 & 0.6143 & 0.5474 & 9.8570 & 385.162 \\
        CDL-CD-0.7 & 0.6825 & 0.5971 & 0.5316 & 9.8570 & 385.162 \\
        CDL-CD-0.6 & 0.6759 & 0.5872 & 0.5223 & 9.8570 & 385.162 \\
        \bottomrule
    \end{tabular}
    \caption{Comparison of performance metrics at different confidence thresholds for the CNNTransroemer-CD-one on the DREAMT dataset.}
    \label{tab:performance_of_different_confidence}
\end{table}

\begin{table}[ht]
    \centering
    \begin{tabular}{lcccccc}
        \toprule
        Method (\(\theta\)) & Accuracy & F1 Score & Cohen's Kappa & FLOPs (M) & Parameters (K) \\
        \midrule
        CDL-Deeponly & 0.6965 & 0.6207 & 0.5450 & 9.8566 & 384.837\\
        CDL-CD-2 & 0.7029 & 0.6309 & 0.5587 & 9.8566 & 384.837\\
        CDL-CD-1 & 0.6944 & 0.6276 & 0.5504 & 9.8566 & 384.837\\
        CDL-CD-0.9 & 0.6901 & 0.6034 & 0.5416 & 9.8573 & 385.482 \\
        CDL-CD-0.8 & 0.6721 & 0.5828 & 0.5191 & 9.8573 & 385.482 \\
        CDL-CD-0.7 & 0.6557 & 0.5671 & 0.5009 & 9.8573 & 385.482 \\
        CDL-CD-0.6 & 0.6473 & 0.5538 & 0.4859 & 9.8573 & 385.482 \\
        \bottomrule
    \end{tabular}
    \caption{Comparison of performance metrics at different confidence thresholds for the CNNTransroemer-CD-two on the DREAMT dataset.}
    \label{tab:performance_of_different_confidence}
\end{table}

\begin{table}[ht]
    \centering
    \begin{tabular}{lcccccc}
        \toprule
        Method (\(\theta\)) & Accuracy & F1 Score & Cohen's Kappa & FLOPs (M) & Parameters (K) \\
        \midrule
        CDL-Deeponly & 0.6965 & 0.6207 & 0.5450 & 9.8566 & 384.837\\
        CDL-CD-2 & 0.7042 & 0.6329 & 0.5605 & 9.8566 & 384.837\\
        CDL-CD-1 & 0.6952 & 0.6261 & 0.5519 & 9.8566 & 384.837\\
        CDL-CD-0.9 & 0.6850 & 0.6080 & 0.5376 & 9.8573 & 385.482 \\
        CDL-CD-0.8 & 0.6617 & 0.5822 & 0.5118 & 9.8573 & 385.482 \\
        CDL-CD-0.7 & 0.6535 & 0.5626 & 0.4975 & 9.8573 & 385.482 \\
        CDL-CD-0.6 & 0.6468 & 0.5517 & 0.4854 & 9.8573 & 385.482 \\
        \bottomrule
    \end{tabular}
    \caption{Comparison of performance metrics at different confidence thresholds for the CNNTransroemer-CD-three on the DREAMT dataset.}
    \label{tab:performance_of_different_confidence}
\end{table}

\begin{table}[ht]
    \centering
    \begin{tabular}{lcccccc}
        \toprule
        Method (\(\theta\)) & Accuracy & F1 Score & Cohen's Kappa & FLOPs (M) & Parameters (K) \\
        \midrule
        CDL-Deeponly & 0.7944 & 0.7220 & 0.7030 & 9.8566 & 384.837\\
        CDL-CD-2 & 0.7944 & 0.7265 & 0.7086 & 9.8566 & 384.837\\
        CDL-CD-1 & 0.7980 & 0.7271 & 0.7091 & 9.8566 & 384.837\\
        CDL-CD-0.9 & 0.7978 & 0.7166 & 0.6984 & 9.8573 & 385.482 \\
        CDL-CD-0.8 & 0.7889 & 0.7097 & 0.6955 & 9.8573 & 385.482 \\
        CDL-CD-0.7 & 0.7821 & 0.7083 & 0.6876 & 9.8573 & 385.482 \\
        CDL-CD-0.6 & 0.7589 & 0.6808 & 0.6547 & 9.8573 & 385.482 \\
        \bottomrule
    \end{tabular}
    \caption{Comparison of performance metrics at different confidence thresholds for the CNNTransroemer-CD-three on the DREAMT dataset.}
    \label{tab:performance_of_different_confidence}
\end{table}
\chapter{Conclusion}

\hspace{2em}This thesis presents a novel approach to power-efficient sleep stage classification designed for consumer wearable devices. The primary focus has been on developing a deep learning model that balances computational efficiency with classification accuracy, making it suitable for deployment in resource-constrained environments.

The research explored the integration of 3-axis motion data, heart rate, and temporal information relative to light-off as input features. A key contribution of this work is the hierarchical model architecture, which employs a combination of simple and deep classifiers. This architecture allows the model to handle simple classification tasks with minimal computational resources, while reserving more complex feature extraction and classification for a deep classifier. By introducing a confidence-based decision mechanism, the model can dynamically determine whether further processing is necessary, thus optimizing both performance and efficiency.

The implementation of confidence thresholds has proven to be effective in managing the trade-off between computational cost and classification performance. A threshold of 0.8 was identified as providing an optimal balance, improving F1 score and Cohen's kappa while maintaining a low computational footprint. Additionally, the use of multiple simple classifiers at various stages of the CNNTransformer architecture demonstrated the benefits of hierarchical design in enhancing model performance for more challenging tasks.

The performance comparison between the proposed method and the CNNTransformer baseline highlighted significant reductions in computational cost—by approximately 80\%—and in the number of parameters—by nearly 90\%—on the SleepAccel dataset. While there was a slight trade-off in accuracy, the hierarchical design and confidence-based mechanism led to notable improvements in precision, recall, and overall classification reliability.

In conclusion, the methods developed in this thesis offer a viable solution for power-efficient sleep stage classification on consumer wearable devices. By leveraging hierarchical structures, parameter sharing, and confidence-based decision making, the proposed approach effectively balances the demands of computational efficiency and classification accuracy, paving the way for more practical and scalable sleep monitoring solutions.

%\bibliographystyle{IEEEtran}  % 使用 IEEE Trans 期刊格式
%\bibliography{texFolder/reference}

% back pages 後頁
% 包括參考文獻、附錄、自傳
% 實際內容由
%    my_bib.bib, my_appendix.tex, my_vita.tex
% 決定
% yzu_backpages.tex 此檔只提供整體架構的定義，不需更動
% 在撰寫各章草稿時，可以把此部份「關掉」，以節省無謂的編譯時間。
% %
% this file is encoded in utf-8
% v2.02 (Sep. 12, 2012)

%%% 參考文獻
\newpage
\phantomsection % for hyperref to register this
\addcontentsline{toc}{chapter}{\nameRef}
\renewcommand{\bibname}{\protect\makebox[5cm][s]{\nameRef}}
%  \makebox{} is fragile; need protect
%\bibliographystyle{IEEEtran}  % 使用 IEEE Trans 期刊格式
%\bibliography{reference}


%%% 附錄
%%
% this file is encoded in utf-8
% v2.02 (Sep. 12, 2012)
%%% 每一個附錄 (附錄甲、附錄乙、...) 都要複製此段附錄編排碼做為起頭
%%% 附錄編排碼 begin >>>
\newpage 
\chapter*{附錄 A:　　MATLAB / Octave 程式列表} % 修改附錄編號與你的附錄名
\phantomsection % for hyperref to register this
\addcontentsline{toc}{chapter}{附錄 A: MATLAB / Octave 程式列表} %建議此內容應與上行相同
%\setcounter{chapter}{0}  %如果用的是 TeXLive2007 則打開此行以避免錯誤 
\setcounter{equation}{0} 
\setcounter{figure}{0} 
\setcounter{footnote}{0} 
\setcounter{section}{0} 
\setcounter{subsection}{0}
\setcounter{subsubsection}{0}
\setcounter{table}{0} 
\renewcommand{\thechapter}{A} % 如果是附錄 B，則內容應為{B}
%%% <<< 附錄編排碼 end

% 附錄內容開始
%%% 納入程式源碼
\lstinputlisting[caption={matlab 程式碼列表範例},
label=lst:matlab:example,
numbers=left,
firstnumber=1,
frame=ltrb, % single lines for left, top, right, bottom; LTRB for double lines 
escapeinside={$$}, %如要在列表裡顯示特殊字元/排版效果，要把該文字串用 $$ 包夾住 (適合 C 程式碼)(原預設為 <>)
]
{example_prog_list.m}

\begin{equation}\sum_{k=1}^{n} k = \frac{n(n+1)}{2}\end{equation}

%%% 如果有附錄B、C、...，則在此繼續加上「附錄編排」碼
% 每一個附錄會自動以新頁開始

%%% 自傳
%\newpage
%\chapter*{\protect\makebox[5cm][s]{\nameVita}} % \makebox{} is fragile; need protect
%\phantomsection % for hyperref to register this
%\addcontentsline{toc}{chapter}{\nameVita}
%\input{my_vita.tex}





\clearpage % to make sure all CJK characters are processed
\end{CJK*}  %%% ZZZ %%%
\printbibliography
\end{document} 
 
